%        1         2         3         4         5         6         7         8         9         0         1         2
%23456789012345678901234567890123456789012345678901234567890123456789012345678901234567890123456789012345678901234567890
\documentclass[12pt]{MIA-USA}

\usepackage{layout}
\usepackage{subcaption}
\usepackage{paralist}
\usepackage{breakcites}

% Paquetes y macros que requiere el contenido convertido desde Markdown (Pandoc)
% ============================================================================
% pandoc-preamble.tex
% Paquetes y macros que requiere el LaTeX generado por Pandoc a partir de los
% archivos Markdown del informe (tablas booktabs/longtable, bloques de código
% con resaltado de sintaxis, listas compactas, etc.).
%
% La clase MIA-USA.cls ya carga: graphicx, xcolor[table], longtable, array,
% colortbl, tabularx, amsmath, tikz, fancyhdr, multirow, multicol, tcolorbox.
% Aquí solo se añade lo que falta. Se usa pdfLaTeX (no se importan fuentes
% Times New Roman/fontspec del metadata.yaml, que eran para XeLaTeX).
% ============================================================================

% --- Tablas generadas por Pandoc -------------------------------------------
\usepackage{booktabs}        % \toprule \midrule \bottomrule
\usepackage{calc}            % cálculos de ancho de columnas de Pandoc
\providecommand{\real}[1]{#1}

% Pandoc define columnas de tablas con tipos personalizados; estos comandos
% reproducen su plantilla por defecto para que las tablas compilen.
\newlength{\cslhangindent}
\setlength{\cslhangindent}{1.5em}

% Pandoc envuelve las tablas sin título en {\def\LTcaptype{none} ...}. Como la
% plantilla carga el paquete caption (vía subcaption), este intenta usar un
% contador 'none' que no existe. Lo definimos para evitar el error fatal.
\makeatletter
\@ifundefined{c@none}{\newcounter{none}}{}
\makeatother

% --- Tablas anchas: que quepan en la página --------------------------------
\usepackage{etoolbox}
\AtBeginEnvironment{longtable}{\footnotesize\setlength{\tabcolsep}{3pt}\renewcommand{\arraystretch}{0.95}}

% Tablas en orientación horizontal cuando el ancho lo requiere
\usepackage{pdflscape}

% --- Bloques de código con resaltado (salida de Pandoc) --------------------
\usepackage{fvextra}
% Ajuste automático de línea para que el código no se desborde del margen
\DefineVerbatimEnvironment{Highlighting}{Verbatim}{breaklines,breakanywhere,commandchars=\\\{\}}
\newcommand{\VerbBar}{|}
\newcommand{\VERB}{\Verb[commandchars=\\\{\}]}
\newenvironment{Shaded}{}{}
\newcommand{\AlertTok}[1]{\textcolor[rgb]{1.00,0.00,0.00}{\textbf{#1}}}
\newcommand{\AnnotationTok}[1]{\textcolor[rgb]{0.38,0.63,0.69}{\textbf{\textit{#1}}}}
\newcommand{\AttributeTok}[1]{\textcolor[rgb]{0.49,0.56,0.16}{#1}}
\newcommand{\BaseNTok}[1]{\textcolor[rgb]{0.25,0.63,0.44}{#1}}
\newcommand{\BuiltInTok}[1]{\textcolor[rgb]{0.00,0.50,0.00}{#1}}
\newcommand{\CharTok}[1]{\textcolor[rgb]{0.25,0.44,0.63}{#1}}
\newcommand{\CommentTok}[1]{\textcolor[rgb]{0.38,0.63,0.69}{\textit{#1}}}
\newcommand{\CommentVarTok}[1]{\textcolor[rgb]{0.38,0.63,0.69}{\textbf{\textit{#1}}}}
\newcommand{\ConstantTok}[1]{\textcolor[rgb]{0.53,0.00,0.00}{#1}}
\newcommand{\ControlFlowTok}[1]{\textcolor[rgb]{0.00,0.44,0.13}{\textbf{#1}}}
\newcommand{\DataTypeTok}[1]{\textcolor[rgb]{0.56,0.13,0.00}{#1}}
\newcommand{\DecValTok}[1]{\textcolor[rgb]{0.25,0.63,0.44}{#1}}
\newcommand{\DocumentationTok}[1]{\textcolor[rgb]{0.73,0.13,0.13}{\textit{#1}}}
\newcommand{\ErrorTok}[1]{\textcolor[rgb]{1.00,0.00,0.00}{\textbf{#1}}}
\newcommand{\ExtensionTok}[1]{#1}
\newcommand{\FloatTok}[1]{\textcolor[rgb]{0.25,0.63,0.44}{#1}}
\newcommand{\FunctionTok}[1]{\textcolor[rgb]{0.02,0.16,0.49}{#1}}
\newcommand{\ImportTok}[1]{\textcolor[rgb]{0.00,0.50,0.00}{\textbf{#1}}}
\newcommand{\InformationTok}[1]{\textcolor[rgb]{0.38,0.63,0.69}{\textbf{\textit{#1}}}}
\newcommand{\KeywordTok}[1]{\textcolor[rgb]{0.00,0.44,0.13}{\textbf{#1}}}
\newcommand{\NormalTok}[1]{#1}
\newcommand{\OperatorTok}[1]{\textcolor[rgb]{0.40,0.40,0.40}{#1}}
\newcommand{\OtherTok}[1]{\textcolor[rgb]{0.00,0.44,0.13}{#1}}
\newcommand{\PreprocessorTok}[1]{\textcolor[rgb]{0.74,0.48,0.00}{#1}}
\newcommand{\RegionMarkerTok}[1]{#1}
\newcommand{\SpecialCharTok}[1]{\textcolor[rgb]{0.25,0.44,0.63}{#1}}
\newcommand{\SpecialStringTok}[1]{\textcolor[rgb]{0.73,0.40,0.53}{#1}}
\newcommand{\StringTok}[1]{\textcolor[rgb]{0.25,0.44,0.63}{#1}}
\newcommand{\VariableTok}[1]{\textcolor[rgb]{0.10,0.09,0.49}{#1}}
\newcommand{\VerbatimStringTok}[1]{\textcolor[rgb]{0.25,0.44,0.63}{#1}}
\newcommand{\WarningTok}[1]{\textcolor[rgb]{0.38,0.63,0.69}{\textbf{\textit{#1}}}}

% --- Listas compactas y varios de Pandoc -----------------------------------
\providecommand{\tightlist}{%
  \setlength{\itemsep}{0pt}\setlength{\parskip}{0pt}}
\setlength{\emergencystretch}{3em} % evita líneas desbordadas

% Caracteres de dibujo de cajas (─ │ └ ├ ┌ ┐ …) usados en los árboles de
% directorios dentro de los bloques de código. pmboxdraw los define para
% pdfLaTeX+inputenc.
\usepackage{pmboxdraw}

% Otros caracteres Unicode que la fuente base / inputenc no traen.
\usepackage{newunicodechar}
\newunicodechar{⋮}{\ensuremath{\vdots}}        % U+22EE elipsis vertical
\newunicodechar{→}{\ensuremath{\rightarrow}}   % U+2192 flecha derecha
\newunicodechar{←}{\ensuremath{\leftarrow}}    % U+2190
\newunicodechar{↔}{\ensuremath{\leftrightarrow}}
\newunicodechar{−}{\ensuremath{-}}             % U+2212 signo menos
\newunicodechar{≤}{\ensuremath{\leq}}          % U+2264
\newunicodechar{≥}{\ensuremath{\geq}}          % U+2265
\newunicodechar{×}{\ensuremath{\times}}        % U+00D7 (por si acaso)
\newunicodechar{►}{\ensuremath{\blacktriangleright}} % U+25BA
\newunicodechar{✅}{\checkmark}                 % U+2705
\newunicodechar{❌}{\ensuremath{\times}}        % U+274C
\newunicodechar{⁴}{\textsuperscript{4}}        % U+2074 superíndices (notas)
\newunicodechar{⁵}{\textsuperscript{5}}
\newunicodechar{⁶}{\textsuperscript{6}}
\newunicodechar{⁷}{\textsuperscript{7}}
\newunicodechar{⁸}{\textsuperscript{8}}
\newunicodechar{⁹}{\textsuperscript{9}}
\usepackage{amssymb}                            % \checkmark, \blacktriangleright

% Pandoc puede emitir \passthrough en código en línea
\providecommand{\passthrough}[1]{#1}

% Enlaces de imágenes/figuras de Pandoc dentro de \includegraphics
\makeatletter
\providecommand{\pandocbounded}[1]{#1}
\makeatother


\usepackage[pagebackref=true,breaklinks=true]{hyperref}
\hypersetup{
	pdftitle={OrbitEngine: Anexos},
	pdfauthor={MIA-D.Martinez},
	pdfsubject={},
	pdfcreator={DarMart with Latex},
	colorlinks=true,
	linkcolor=GrisUno,
	citecolor=AzulInstitucional,
	filecolor=AzulClaro,
	urlcolor=AzulInstitucional,
	linktoc=all
}

\graphicspath{{Images/}}


\title{Orbit Engine: Plataforma SaaS para la Gesti\'on Integral de Procesos Internos en Peque\~nas y Medianas Empresas}
\author{Rinc\'on Suarez Fabi\'an \\
Rodr\'iguez Forero Nicol\'as \\
Velasco Gonz\'alez Daniel}
\documenttype{Anexos}

\advisor{Juan Pablo Ospina}

\coadvisor{Camilo Rodr\'iguez}


% Aqui comienza el documento
% - -- --- ----- ------- ------- ----- --- -- -
\begin{document}

	\maketitle

    % Crea la tabla de contenidos a partir de la estructura del documento
    \tableofcontents

    % Lista de tablas (los anexos no contienen figuras).
    % tocbibind (cargado por la clase) ya la añade automáticamente al índice.
    \cleardoublepage
    \listoftables

    % - -- --- ----- Anexos (A, B, C) ----- --- -- -
    % Convertidos automáticamente desde docs/informe-final/anexo-*.md con Pandoc.
    % \appendix hace que los capítulos se numeren A, B, C.

    \appendix

    \chapter{Manual de Usuario}\label{manual-de-usuario}

\textbf{OrbitEngine} --- Plataforma SaaS para la Gestión de Procesos Internos en Pymes\\
Versión: 1.0 \textbar{} Abril 2026\\
Universidad Sergio Arboleda --- Semillero de Software como Innovación

\begin{center}\rule{0.5\linewidth}{0.5pt}\end{center}

\section{Introducción}\label{introducciuxf3n}

Este manual describe el uso de OrbitEngine para los usuarios finales de la plataforma. OrbitEngine es una aplicación web multi-tenant que permite a pequeñas y medianas empresas gestionar en un solo lugar sus operaciones de \textbf{inventario}, \textbf{ventas} y \textbf{clientes}, con un panel de indicadores clave de desempeño (KPIs) accesible en tiempo real.

El sistema opera completamente desde el navegador web, sin necesidad de instalar software adicional. Cada empresa (organización) tiene su propio espacio de trabajo completamente aislado de las demás.

\section{Acceso a la Plataforma}\label{acceso-a-la-plataforma}

\subsection{Página de inicio}\label{puxe1gina-de-inicio}

Al ingresar a la URL de OrbitEngine, el usuario accede a la \textbf{página de presentación} (landing page) que muestra las características de la plataforma. Desde esta página se puede navegar a las opciones de inicio de sesión o registro.

\subsection{Registro de organización}\label{registro-de-organizaciuxf3n}

Para crear una nueva empresa en la plataforma:

\begin{enumerate}
\def\labelenumi{\arabic{enumi}.}
\tightlist
\item
  Hacer clic en \textbf{``Comenzar gratis''} o en el enlace de registro en la barra de navegación.
\item
  En la pantalla \texttt{/signup-org}, completar el formulario con:

  \begin{itemize}
  \tightlist
  \item
    \textbf{Nombre de la organización} --- nombre visible de la empresa.
  \item
    \textbf{Slug} --- identificador único en minúsculas (ej. \texttt{mi-empresa}). Solo letras, números y guiones; entre 3 y 50 caracteres.
  \item
    \textbf{Descripción} \emph{(opcional)} --- breve descripción del negocio.
  \item
    \textbf{Nombre}, \textbf{apellido} y \textbf{correo electrónico} del administrador.
  \item
    \textbf{Contraseña} de al menos 8 caracteres.
  \end{itemize}
\item
  Hacer clic en \textbf{``Crear organización''}.
\item
  El sistema crea la organización y la cuenta administrador automáticamente y redirige al panel de control.
\end{enumerate}

\subsection{Invitación de usuarios}\label{invitaciuxf3n-de-usuarios}

Los usuarios adicionales \textbf{no se registran por su cuenta}: son creados por el administrador desde el módulo \textbf{Admin}. Una vez creada la cuenta, el usuario recibe las credenciales y puede ingresar con ellas.

\subsection{Inicio de sesión}\label{inicio-de-sesiuxf3n}

\begin{enumerate}
\def\labelenumi{\arabic{enumi}.}
\tightlist
\item
  Ir a la pantalla \texttt{/login}.
\item
  Ingresar el \textbf{correo electrónico} y la \textbf{contraseña}.
\item
  Hacer clic en \textbf{``Ingresar''}.
\item
  Si las credenciales son correctas, el sistema redirige al \textbf{panel de control} (\texttt{/dashboard}).
\end{enumerate}

\subsection{Recuperación de contraseña}\label{recuperaciuxf3n-de-contraseuxf1a}

Si olvidó su contraseña:

\begin{enumerate}
\def\labelenumi{\arabic{enumi}.}
\tightlist
\item
  En la pantalla de inicio de sesión, hacer clic en \textbf{``¿Olvidaste tu contraseña?''}.
\item
  Ingresar el correo electrónico asociado a la cuenta.
\item
  Revisar la bandeja de entrada y hacer clic en el enlace recibido.
\item
  En la pantalla \texttt{/reset-password}, ingresar y confirmar la nueva contraseña.
\item
  Hacer clic en \textbf{``Restablecer contraseña''}.
\end{enumerate}

\subsection{Cierre de sesión}\label{cierre-de-sesiuxf3n}

Para cerrar sesión, hacer clic en el \textbf{ícono de usuario} en la barra lateral izquierda y seleccionar \textbf{``Cerrar sesión''}. La sesión expira automáticamente después de 8 días de inactividad.

\section{Navegación General}\label{navegaciuxf3n-general}

Una vez dentro del sistema, la interfaz se divide en:

\begin{itemize}
\tightlist
\item
  \textbf{Barra lateral izquierda} --- menú principal de navegación con acceso a todos los módulos.
\item
  \textbf{Área de contenido principal} --- donde se muestra y gestiona la información de cada módulo.
\item
  \textbf{Panel de usuario} (parte inferior de la barra lateral) --- acceso a configuración personal y cierre de sesión.
\end{itemize}

Los módulos disponibles son:

{\def\LTcaptype{none} % do not increment counter
\begin{longtable}[]{@{}lll@{}}
\toprule\noalign{}
Módulo & Ruta & Acceso \\
\midrule\noalign{}
\endhead
\bottomrule\noalign{}
\endlastfoot
Dashboard & \texttt{/dashboard} & Todos los usuarios \\
Inventario & \texttt{/dashboard/inventory} & Todos los usuarios \\
Ventas & \texttt{/dashboard/sales} & Todos los usuarios \\
Clientes & \texttt{/dashboard/customers} & Todos los usuarios \\
Administración & \texttt{/dashboard/admin} & Solo Administradores \\
Configuración & \texttt{/dashboard/settings} & Todos los usuarios \\
\end{longtable}
}

\section{Dashboard --- Panel de Indicadores}\label{dashboard-panel-de-indicadores}

El dashboard muestra un resumen del estado del negocio en tiempo real para la organización activa.

\textbf{Indicadores disponibles:}

\begin{itemize}
\tightlist
\item
  \textbf{Ventas del día} --- total en dinero de las ventas registradas hoy.
\item
  \textbf{Ventas del mes} --- total acumulado del mes en curso.
\item
  \textbf{Ticket promedio} --- valor promedio por venta en el período.
\item
  \textbf{Productos con stock bajo} --- conteo de productos que están por debajo del mínimo configurado.
\item
  \textbf{Top productos} --- listado de los productos más vendidos.
\item
  \textbf{Ventas por día} --- gráfica de barras con la evolución diaria de ventas.
\end{itemize}

El dashboard se actualiza automáticamente con cada consulta. No requiere acción adicional del usuario.

\section{Módulo de Inventario}\label{muxf3dulo-de-inventario}

El módulo de inventario permite gestionar los productos y categorías de la empresa, así como llevar el control detallado de los movimientos de stock.

\subsection{Categorías}\label{categoruxedas}

Las categorías permiten organizar el inventario en grupos lógicos. Admiten \textbf{jerarquía} (una categoría puede tener subcategorías).

\textbf{Crear una categoría:}

\begin{enumerate}
\def\labelenumi{\arabic{enumi}.}
\tightlist
\item
  En el módulo Inventario, ir a la pestaña \textbf{Categorías}.
\item
  Hacer clic en \textbf{``Agregar categoría''}.
\item
  Completar el \textbf{nombre} \emph{(requerido)} y opcionalmente una \textbf{descripción} y una \textbf{categoría padre}.
\item
  Guardar.
\end{enumerate}

\textbf{Editar o desactivar una categoría:}

\begin{itemize}
\tightlist
\item
  Hacer clic en el menú de acciones (⋮) de la fila y seleccionar \textbf{``Editar''} o \textbf{``Desactivar''}.
\end{itemize}

\subsection{Productos}\label{productos}

\textbf{Crear un producto:}

\begin{enumerate}
\def\labelenumi{\arabic{enumi}.}
\tightlist
\item
  En el módulo Inventario, ir a la pestaña \textbf{Productos}.
\item
  Hacer clic en \textbf{``Agregar producto''}.
\item
  Completar el formulario:

  \begin{itemize}
  \tightlist
  \item
    \textbf{Nombre} \emph{(requerido)}
  \item
    \textbf{SKU} --- código único del producto \emph{(requerido)}
  \item
    \textbf{Categoría} \emph{(opcional)}
  \item
    \textbf{Precio de costo} y \textbf{precio de venta}
  \item
    \textbf{Stock actual}, \textbf{stock mínimo} y \textbf{stock máximo} \emph{(opcional)}
  \item
    \textbf{Unidad de medida} (ej. \texttt{unit}, \texttt{kg}, \texttt{lt})
  \item
    \textbf{Código de barras} \emph{(opcional)}
  \item
    \textbf{Descripción} e \textbf{imagen} \emph{(opcionales)}
  \end{itemize}
\item
  Guardar.
\end{enumerate}

\textbf{Editar un producto:}

\begin{itemize}
\tightlist
\item
  Hacer clic en el menú de acciones (⋮) → \textbf{``Editar''}. Todos los campos son modificables excepto el SKU si ya tiene movimientos.
\end{itemize}

\textbf{Desactivar un producto:}

\begin{itemize}
\tightlist
\item
  Hacer clic en el menú de acciones (⋮) → \textbf{``Desactivar''}. El producto deja de aparecer en los formularios de venta pero su historial se conserva.
\end{itemize}

\textbf{Alerta de stock bajo:}

\begin{itemize}
\tightlist
\item
  Los productos con \texttt{stock\_quantity\ ≤\ stock\_min} aparecen resaltados en la tabla y se contabilizan en el indicador del dashboard.
\end{itemize}

\subsection{Ajuste de stock manual}\label{ajuste-de-stock-manual}

Para corregir el inventario por diferencias de conteo, mermas u otros motivos:

\begin{enumerate}
\def\labelenumi{\arabic{enumi}.}
\tightlist
\item
  En la tabla de productos, seleccionar el producto y hacer clic en \textbf{``Ajustar stock''}.
\item
  Ingresar la \textbf{cantidad} (positiva para agregar, negativa para restar).
\item
  Ingresar un \textbf{motivo} del ajuste \emph{(requerido)}.
\item
  Confirmar.
\end{enumerate}

El ajuste queda registrado en el historial de movimientos con el tipo \texttt{adjustment}.

\subsection{Historial de movimientos}\label{historial-de-movimientos}

El historial registra todos los cambios de stock de cada producto:

{\def\LTcaptype{none} % do not increment counter
\begin{longtable}[]{@{}ll@{}}
\toprule\noalign{}
Tipo de movimiento & Descripción \\
\midrule\noalign{}
\endhead
\bottomrule\noalign{}
\endlastfoot
\texttt{sale} & Salida generada por una venta \\
\texttt{adjustment} & Entrada o salida manual \\
\texttt{return} & Entrada por devolución o cancelación de venta \\
\texttt{purchase} & Entrada por compra \\
\end{longtable}
}

Para ver el historial de un producto:

\begin{enumerate}
\def\labelenumi{\arabic{enumi}.}
\tightlist
\item
  En la tabla de productos, hacer clic en el menú de acciones (⋮) → \textbf{``Ver historial de movimientos''}.
\item
  Se muestra la tabla con fecha, tipo, cantidad, stock anterior y stock resultante, y el usuario que realizó el movimiento.
\end{enumerate}

\section{Módulo de Ventas}\label{muxf3dulo-de-ventas}

El módulo de ventas permite registrar transacciones comerciales, visualizar el historial y cancelar ventas cuando sea necesario.

\subsection{Registrar una venta}\label{registrar-una-venta}

\begin{enumerate}
\def\labelenumi{\arabic{enumi}.}
\tightlist
\item
  En el módulo Ventas, hacer clic en \textbf{``Nueva venta''}.
\item
  Completar el formulario:

  \begin{itemize}
  \tightlist
  \item
    \textbf{Cliente} \emph{(opcional)} --- buscar por nombre o documento.
  \item
    \textbf{Método de pago}: Efectivo, Tarjeta, Transferencia u Otro.
  \item
    \textbf{Descuento} y \textbf{IVA/impuesto} \emph{(opcional, en valores monetarios)}.
  \item
    \textbf{Notas} \emph{(opcional)}.
  \end{itemize}
\item
  En la sección de \textbf{ítems}, agregar los productos de la venta:

  \begin{itemize}
  \tightlist
  \item
    Buscar el producto por nombre o SKU.
  \item
    Ingresar la \textbf{cantidad}.
  \item
    El sistema calcula el precio unitario y el subtotal del ítem automáticamente.
  \end{itemize}
\item
  El sistema calcula en tiempo real: subtotal, descuento, impuesto y \textbf{total}.
\item
  Hacer clic en \textbf{``Registrar venta''}.
\end{enumerate}

Al guardar, el sistema:

\begin{itemize}
\tightlist
\item
  Genera un \textbf{número de factura} único.
\item
  Descuenta las cantidades del stock de cada producto registrado.
\item
  Registra los movimientos de inventario correspondientes (\texttt{sale}).
\item
  Actualiza el total de compras del cliente, si se asoció uno.
\end{itemize}

\subsection{Ver detalle de una venta}\label{ver-detalle-de-una-venta}

En la tabla de ventas, hacer clic en el menú de acciones (⋮) → \textbf{``Ver detalle''}. Se muestra la información completa: número de factura, cliente, fecha, ítems, totales y método de pago.

\subsection{Cancelar una venta}\label{cancelar-una-venta}

Una venta puede cancelarse mientras tenga estado \texttt{completed}:

\begin{enumerate}
\def\labelenumi{\arabic{enumi}.}
\tightlist
\item
  En la tabla de ventas, hacer clic en el menú de acciones (⋮) → \textbf{``Cancelar venta''}.
\item
  Ingresar el \textbf{motivo de cancelación} \emph{(requerido)}.
\item
  Confirmar.
\end{enumerate}

Al cancelar, el sistema:

\begin{itemize}
\tightlist
\item
  Cambia el estado de la venta a \texttt{cancelled}.
\item
  \textbf{Revierte el stock} de cada producto afectado (movimiento tipo \texttt{return}).
\item
  Registra la fecha y usuario de cancelación.
\end{itemize}

\subsection{Filtros y búsqueda}\label{filtros-y-buxfasqueda}

La tabla de ventas permite:

\begin{itemize}
\tightlist
\item
  Buscar por \textbf{número de factura}.
\item
  Filtrar por \textbf{estado} (\texttt{completed}, \texttt{cancelled}) y \textbf{método de pago}.
\item
  Ordenar por cualquier columna.
\end{itemize}

\section{Módulo de Clientes}\label{muxf3dulo-de-clientes}

El módulo permite gestionar el directorio de clientes de la organización y consultar su historial de compras.

\subsection{Registrar un cliente}\label{registrar-un-cliente}

\begin{enumerate}
\def\labelenumi{\arabic{enumi}.}
\tightlist
\item
  En el módulo Clientes, hacer clic en \textbf{``Agregar cliente''}.
\item
  Completar el formulario:

  \begin{itemize}
  \tightlist
  \item
    \textbf{Tipo de documento}: CC, NIT, Pasaporte, Otro \emph{(requerido)}.
  \item
    \textbf{Número de documento} \emph{(requerido, único por organización)}.
  \item
    \textbf{Nombre} y \textbf{apellido} \emph{(requeridos)}.
  \item
    \textbf{Correo electrónico}, \textbf{teléfono}, \textbf{dirección}, \textbf{ciudad} y \textbf{país} \emph{(opcionales)}.
  \item
    \textbf{Notas internas} \emph{(opcionales)}.
  \end{itemize}
\item
  Guardar.
\end{enumerate}

\subsection{Editar un cliente}\label{editar-un-cliente}

\begin{itemize}
\tightlist
\item
  Hacer clic en el menú de acciones (⋮) → \textbf{``Editar''}. Se pueden modificar todos los campos excepto el número de documento.
\end{itemize}

\subsection{Desactivar un cliente}\label{desactivar-un-cliente}

\begin{itemize}
\tightlist
\item
  Hacer clic en el menú de acciones (⋮) → \textbf{``Desactivar''}. El cliente se desactiva pero su historial de compras se conserva.
\end{itemize}

\subsection{Historial de compras del cliente}\label{historial-de-compras-del-cliente}

\begin{itemize}
\tightlist
\item
  Hacer clic en el menú de acciones (⋮) → \textbf{``Ver historial de compras''}.
\item
  Se muestran todas las ventas asociadas al cliente con fecha, número de factura, total y estado.
\item
  Se visualizan también los indicadores: \textbf{total acumulado de compras}, \textbf{número de compras} y \textbf{fecha de última compra}.
\end{itemize}

\section{Módulo de Administración}\label{muxf3dulo-de-administraciuxf3n}

\begin{quote}
Este módulo es exclusivo para usuarios con el rol \textbf{Administrador}.
\end{quote}

El módulo de administración permite gestionar las cuentas de usuario de la organización.

\subsection{Ver usuarios}\label{ver-usuarios}

La tabla muestra todos los usuarios de la organización con su nombre, correo, rol y estado (activo/inactivo).

\subsection{Crear un usuario}\label{crear-un-usuario}

\begin{enumerate}
\def\labelenumi{\arabic{enumi}.}
\tightlist
\item
  Hacer clic en \textbf{``Agregar usuario''}.
\item
  Completar el formulario:

  \begin{itemize}
  \tightlist
  \item
    \textbf{Nombre} y \textbf{apellido}.
  \item
    \textbf{Correo electrónico} --- será el identificador de ingreso.
  \item
    \textbf{Contraseña} temporal.
  \item
    \textbf{Rol} --- seleccionar entre los roles disponibles de la organización.
  \end{itemize}
\item
  Guardar.
\end{enumerate}

El nuevo usuario puede ingresar inmediatamente con las credenciales proporcionadas y cambiar su contraseña desde la configuración.

\subsection{Editar un usuario}\label{editar-un-usuario}

\begin{itemize}
\tightlist
\item
  Hacer clic en el menú de acciones (⋮) → \textbf{``Editar''}. Se pueden cambiar nombre, apellido y rol.
\end{itemize}

\subsection{Eliminar un usuario}\label{eliminar-un-usuario}

\begin{itemize}
\tightlist
\item
  Hacer clic en el menú de acciones (⋮) → \textbf{``Eliminar''}.
\item
  Confirmar la acción en el diálogo de confirmación.
\end{itemize}

\begin{quote}
\textbf{Advertencia:} Esta acción es permanente. El historial de ventas y movimientos del usuario permanece registrado.
\end{quote}

\section{Configuración}\label{configuraciuxf3n}

El módulo de configuración permite al usuario gestionar su cuenta personal y, si es administrador, los datos de la organización.

\subsection{Información personal}\label{informaciuxf3n-personal}

En la pestaña \textbf{``Perfil''}:

\begin{itemize}
\tightlist
\item
  Ver y editar \textbf{nombre} y \textbf{apellido}.
\item
  Ver el \textbf{correo electrónico} asociado (no editable desde aquí).
\end{itemize}

\subsection{Cambio de contraseña}\label{cambio-de-contraseuxf1a}

En la pestaña \textbf{``Seguridad''}:

\begin{enumerate}
\def\labelenumi{\arabic{enumi}.}
\tightlist
\item
  Ingresar la \textbf{contraseña actual}.
\item
  Ingresar la \textbf{nueva contraseña} y confirmarla.
\item
  Hacer clic en \textbf{``Actualizar contraseña''}.
\end{enumerate}

\subsection{Configuración de la organización}\label{configuraciuxf3n-de-la-organizaciuxf3n}

\begin{quote}
Disponible solo para usuarios con rol \textbf{Administrador}.
\end{quote}

En la pestaña \textbf{``Organización''}:

\begin{itemize}
\tightlist
\item
  Editar el \textbf{nombre} de la organización.
\item
  Editar el \textbf{slug} \emph{(identificador único)}.
\item
  Editar la \textbf{descripción} y la \textbf{URL del logo}.
\end{itemize}

\subsection{Apariencia}\label{apariencia}

En la pestaña \textbf{``Apariencia''}:

\begin{itemize}
\tightlist
\item
  Cambiar entre \textbf{modo claro} y \textbf{modo oscuro}.
\end{itemize}

\subsection{Eliminación de cuenta}\label{eliminaciuxf3n-de-cuenta}

En la pestaña \textbf{``Cuenta''}:

\begin{itemize}
\tightlist
\item
  El usuario puede solicitar la \textbf{eliminación de su propia cuenta}.
\item
  Se requiere confirmación escribiendo el correo electrónico.
\end{itemize}

\begin{quote}
\textbf{Nota:} La eliminación de cuenta no elimina el historial de transacciones asociadas.
\end{quote}

\section{Preguntas Frecuentes}\label{preguntas-frecuentes}

\textbf{¿Puedo usar OrbitEngine desde el celular?}\\
Sí. La interfaz está diseñada de manera responsiva y funciona en navegadores móviles modernos (Chrome, Safari). La experiencia óptima es en pantallas de escritorio o tabletas.

\textbf{¿Mis datos están separados de los de otras empresas?}\\
Sí. Cada organización tiene un espacio de trabajo completamente aislado. No es posible ver ni acceder a datos de otras organizaciones.

\textbf{¿Qué pasa si cancelo una venta? ¿El stock vuelve?}\\
Sí. Al cancelar una venta, el sistema revierte automáticamente el stock de todos los productos incluidos en ella.

\textbf{¿Puedo tener múltiples roles para usuarios?}\\
Actualmente cada usuario tiene un único rol. Los roles definen los permisos de acceso a los módulos y acciones disponibles.

\textbf{¿Cómo sé cuándo un producto está por agotarse?}\\
Los productos cuyo stock actual es igual o menor al \textbf{stock mínimo} configurado aparecen resaltados en la tabla de inventario y se contabilizan en el indicador ``Productos con stock bajo'' del dashboard.

\textbf{¿Se puede exportar la información a Excel o PDF?}\\
Los listados de inventario, ventas y clientes pueden exportarse a Excel (.xlsx) desde el módulo correspondiente, conservando los filtros aplicados. La exportación a PDF no está disponible en esta versión.

\textbf{¿Qué navegadores son compatibles?}\\
OrbitEngine es compatible con las versiones recientes de Chrome, Firefox, Edge y Safari. No se garantiza compatibilidad con Internet Explorer.

\emph{Documento generado como parte del proyecto de grado --- Universidad Sergio Arboleda, Semillero de Software como Innovación, Abril 2026.}

    \chapter{Anexo B. Manual de Instalación y Despliegue}\label{anexo-b.-manual-de-instalaciuxf3n-y-despliegue}

\textbf{OrbitEngine}: Plataforma SaaS para la Gestión de Procesos Internos en Pymes\\
Versión: 1.0 \textbar{} Abril 2026\\
Universidad Sergio Arboleda, Semillero de Software como Innovación

\begin{center}\rule{0.5\linewidth}{0.5pt}\end{center}

\section{Visión General}\label{visiuxf3n-general}

OrbitEngine utiliza \textbf{Docker Compose} como orquestador de contenedores en todos los entornos. La arquitectura de despliegue se compone de los siguientes servicios:

{\def\LTcaptype{none} % do not increment counter
\begin{longtable}[]{@{}
  >{\raggedright\arraybackslash}p{(\linewidth - 4\tabcolsep) * \real{0.1529}}
  >{\raggedright\arraybackslash}p{(\linewidth - 4\tabcolsep) * \real{0.2824}}
  >{\raggedright\arraybackslash}p{(\linewidth - 4\tabcolsep) * \real{0.5647}}@{}}
\toprule\noalign{}
\begin{minipage}[b]{\linewidth}\raggedright
Servicio
\end{minipage} & \begin{minipage}[b]{\linewidth}\raggedright
Imagen
\end{minipage} & \begin{minipage}[b]{\linewidth}\raggedright
Descripción
\end{minipage} \\
\midrule\noalign{}
\endhead
\bottomrule\noalign{}
\endlastfoot
\texttt{db} & \texttt{postgres:18} & Base de datos relacional PostgreSQL \\
\texttt{backend} & Custom (FastAPI) & API REST del sistema \\
\texttt{frontend} & Custom (React/Nginx) & Interfaz de usuario compilada \\
\texttt{prestart} & Custom (FastAPI) & Ejecuta migraciones y datos iniciales al iniciar \\
\texttt{adminer} & \texttt{adminer} & Panel web de administración de BD \\
\texttt{proxy} & \texttt{traefik:3.6} & Reverse proxy y gestor de rutas HTTP/HTTPS \\
\texttt{mailcatcher} & \texttt{schickling/mailcatcher} & Servidor SMTP de prueba (solo dev) \\
\end{longtable}
}

El enrutamiento externo se gestiona con \textbf{Traefik}, que expone los servicios bajo subdominios configurados y maneja la terminación TLS automática mediante Let's Encrypt en producción.

\section{Prerrequisitos}\label{prerrequisitos}

\subsection{Herramientas requeridas}\label{herramientas-requeridas}

{\def\LTcaptype{none} % do not increment counter
\begin{longtable}[]{@{}
  >{\raggedright\arraybackslash}p{(\linewidth - 4\tabcolsep) * \real{0.1489}}
  >{\raggedright\arraybackslash}p{(\linewidth - 4\tabcolsep) * \real{0.1489}}
  >{\raggedright\arraybackslash}p{(\linewidth - 4\tabcolsep) * \real{0.7021}}@{}}
\toprule\noalign{}
\begin{minipage}[b]{\linewidth}\raggedright
Herramienta
\end{minipage} & \begin{minipage}[b]{\linewidth}\raggedright
Versión mínima
\end{minipage} & \begin{minipage}[b]{\linewidth}\raggedright
Propósito
\end{minipage} \\
\midrule\noalign{}
\endhead
\bottomrule\noalign{}
\endlastfoot
Docker Engine & 24.x & Contenedores \\
Docker Compose & v2.x & Orquestación de servicios \\
Git & 2.x & Control de versiones \\
\texttt{uv} & 0.4+ & Gestor de dependencias Python (para desarrollo local sin Docker) \\
Bun & 1.x & Runtime JS / gestor de paquetes (para desarrollo local sin Docker) \\
\end{longtable}
}

\subsection{Puertos utilizados}\label{puertos-utilizados}

{\def\LTcaptype{none} % do not increment counter
\begin{longtable}[]{@{}lll@{}}
\toprule\noalign{}
Puerto & Servicio & Entorno \\
\midrule\noalign{}
\endhead
\bottomrule\noalign{}
\endlastfoot
80 & Traefik HTTP & Dev y Producción \\
443 & Traefik HTTPS & Producción \\
8000 & Backend FastAPI & Dev (expuesto directamente) \\
5173 & Frontend (Vite) & Dev \\
5432 & PostgreSQL & Dev \\
8080 & Adminer & Dev \\
8090 & Traefik Dashboard & Dev \\
1025 & Mailcatcher SMTP & Dev \\
1080 & Mailcatcher UI & Dev \\
\end{longtable}
}

\section{Configuración de Variables de Entorno}\label{configuraciuxf3n-de-variables-de-entorno}

Antes de ejecutar cualquier entorno, crear el archivo \texttt{.env} en la raíz del repositorio:

\begin{Shaded}
\begin{Highlighting}[]
\FunctionTok{cp}\NormalTok{ .env.example .env   }\CommentTok{\# si existe el archivo de ejemplo, de lo contrario crearlo manualmente}
\end{Highlighting}
\end{Shaded}

Las variables requeridas son:

\begin{Shaded}
\begin{Highlighting}[]
\CommentTok{\# ─── Proyecto ──────────────────────────────────────────────────────────────────}
\DataTypeTok{PROJECT\_NAME}\OtherTok{=}\StringTok{OrbitEngine}
\DataTypeTok{STACK\_NAME}\OtherTok{=}\StringTok{orbitengine{-}stack}
\DataTypeTok{DOMAIN}\OtherTok{=}\StringTok{localhost                       \# En producción: tu{-}dominio.com}

\CommentTok{\# ─── Backend ───────────────────────────────────────────────────────────────────}
\DataTypeTok{SECRET\_KEY}\OtherTok{=}\StringTok{cambia{-}esto{-}en{-}produccion   \# Clave aleatoria, mínimo 32 chars}
\DataTypeTok{BACKEND\_CORS\_ORIGINS}\OtherTok{=}\StringTok{["http://localhost:5173","http://localhost"]}
\DataTypeTok{FRONTEND\_HOST}\OtherTok{=}\StringTok{http://localhost:5173    \# En producción: https://tu{-}dominio.com}

\CommentTok{\# ─── Base de datos ─────────────────────────────────────────────────────────────}
\DataTypeTok{POSTGRES\_SERVER}\OtherTok{=}\StringTok{db}
\DataTypeTok{POSTGRES\_PORT}\OtherTok{=}\DecValTok{5432}
\DataTypeTok{POSTGRES\_USER}\OtherTok{=}\StringTok{postgres}
\DataTypeTok{POSTGRES\_PASSWORD}\OtherTok{=}\StringTok{cambia{-}esto}
\DataTypeTok{POSTGRES\_DB}\OtherTok{=}\StringTok{app}

\CommentTok{\# ─── Superusuario inicial ──────────────────────────────────────────────────────}
\DataTypeTok{FIRST\_SUPERUSER}\OtherTok{=}\StringTok{admin@ejemplo.com}
\DataTypeTok{FIRST\_SUPERUSER\_PASSWORD}\OtherTok{=}\StringTok{cambia{-}esto}

\CommentTok{\# ─── Email ─────────────────────────────────────────────────────────────────────}
\DataTypeTok{SMTP\_HOST}\OtherTok{=}\StringTok{mailcatcher                  \# En producción: smtp.tu{-}proveedor.com}
\DataTypeTok{SMTP\_PORT}\OtherTok{=}\StringTok{1025                         \# En producción: 587}
\DataTypeTok{SMTP\_TLS}\OtherTok{=}\StringTok{false                         \# En producción: true}
\DataTypeTok{EMAILS\_FROM\_EMAIL}\OtherTok{=}\StringTok{noreply@tu{-}dominio.com}

\CommentTok{\# ─── Imágenes Docker ──────────────────────────────────────────────────────────}
\DataTypeTok{DOCKER\_IMAGE\_BACKEND}\OtherTok{=}\StringTok{backend}
\DataTypeTok{DOCKER\_IMAGE\_FRONTEND}\OtherTok{=}\StringTok{frontend}
\DataTypeTok{ENVIRONMENT}\OtherTok{=}\StringTok{local                      \# local | staging | production}
\end{Highlighting}
\end{Shaded}

\begin{quote}
\textbf{Producción:} Cambiar todas las claves y contraseñas por valores aleatorios seguros. Se recomienda gestionar las variables de entorno directamente en el panel de Railway (backend) y en el panel de Vercel (frontend), evitando archivos \texttt{.env} en producción.
\end{quote}

\section{Entorno de Desarrollo Local}\label{entorno-de-desarrollo-local}

El entorno de desarrollo combina \texttt{compose.yml} con \texttt{compose.override.yml}, que añade los servicios de desarrollo (Traefik local, Mailcatcher) y sobrescribe las configuraciones para recarga en caliente.

\subsection{Iniciar el entorno}\label{iniciar-el-entorno}

\begin{Shaded}
\begin{Highlighting}[]
\CommentTok{\# Desde la raíz del repositorio}
\ExtensionTok{docker}\NormalTok{ compose watch}
\end{Highlighting}
\end{Shaded}

Este comando:

\begin{enumerate}
\def\labelenumi{\arabic{enumi}.}
\tightlist
\item
  Construye las imágenes de backend y frontend si no existen.
\item
  Levanta todos los servicios en orden de dependencias.
\item
  Activa la recarga automática del backend al detectar cambios en \texttt{./backend}.
\item
  El servicio \texttt{prestart} ejecuta las migraciones de Alembic y los datos iniciales antes de que el backend acepte peticiones.
\end{enumerate}

\subsection{URLs del entorno local}\label{urls-del-entorno-local}

{\def\LTcaptype{none} % do not increment counter
\begin{longtable}[]{@{}ll@{}}
\toprule\noalign{}
Servicio & URL \\
\midrule\noalign{}
\endhead
\bottomrule\noalign{}
\endlastfoot
Frontend & http://localhost:5173 \\
Backend API & http://localhost:8000 \\
Documentación interactiva (Swagger) & http://localhost:8000/docs \\
ReDoc & http://localhost:8000/redoc \\
Adminer & http://localhost:8080 \\
Traefik Dashboard & http://localhost:8090 \\
Mailcatcher UI & http://localhost:1080 \\
\end{longtable}
}

\subsection{Detener el entorno}\label{detener-el-entorno}

\begin{Shaded}
\begin{Highlighting}[]
\CommentTok{\# Detener y eliminar contenedores (datos persisten en volumen)}
\ExtensionTok{docker}\NormalTok{ compose down}

\CommentTok{\# Detener y eliminar contenedores + volúmenes (elimina datos de BD)}
\ExtensionTok{docker}\NormalTok{ compose down }\AttributeTok{{-}v}
\end{Highlighting}
\end{Shaded}

\subsection{Ver logs}\label{ver-logs}

\begin{Shaded}
\begin{Highlighting}[]
\ExtensionTok{docker}\NormalTok{ compose logs }\AttributeTok{{-}f}\NormalTok{ backend    }\CommentTok{\# Logs del backend en tiempo real}
\ExtensionTok{docker}\NormalTok{ compose logs }\AttributeTok{{-}f}\NormalTok{ frontend   }\CommentTok{\# Logs del frontend}
\ExtensionTok{docker}\NormalTok{ compose logs }\AttributeTok{{-}f}\NormalTok{ db         }\CommentTok{\# Logs de PostgreSQL}
\end{Highlighting}
\end{Shaded}

\section{Migraciones de Base de Datos}\label{migraciones-de-base-de-datos}

Las migraciones se gestionan con \textbf{Alembic} y se ejecutan automáticamente al iniciar el entorno (servicio \texttt{prestart}). Para operaciones manuales:

\subsection{Aplicar migraciones pendientes}\label{aplicar-migraciones-pendientes}

\begin{Shaded}
\begin{Highlighting}[]
\ExtensionTok{docker}\NormalTok{ compose exec backend alembic upgrade head}
\end{Highlighting}
\end{Shaded}

\subsection{Crear una nueva migración}\label{crear-una-nueva-migraciuxf3n}

Después de modificar los modelos en \texttt{backend/app/models.py}:

\begin{Shaded}
\begin{Highlighting}[]
\ExtensionTok{docker}\NormalTok{ compose exec backend alembic revision }\AttributeTok{{-}{-}autogenerate} \AttributeTok{{-}m} \StringTok{"descripcion\_del\_cambio"}
\end{Highlighting}
\end{Shaded}

Revisar el archivo generado en \texttt{backend/app/alembic/versions/} antes de aplicarlo.

\subsection{Ver historial de migraciones}\label{ver-historial-de-migraciones}

\begin{Shaded}
\begin{Highlighting}[]
\ExtensionTok{docker}\NormalTok{ compose exec backend alembic history }\AttributeTok{{-}{-}verbose}
\end{Highlighting}
\end{Shaded}

\subsection{Revertir la última migración}\label{revertir-la-uxfaltima-migraciuxf3n}

\begin{Shaded}
\begin{Highlighting}[]
\ExtensionTok{docker}\NormalTok{ compose exec backend alembic downgrade }\AttributeTok{{-}1}
\end{Highlighting}
\end{Shaded}

\subsection{Migraciones del proyecto}\label{migraciones-del-proyecto}

Las migraciones actuales del esquema son:

{\def\LTcaptype{none} % do not increment counter
\begin{longtable}[]{@{}
  >{\raggedright\arraybackslash}p{(\linewidth - 2\tabcolsep) * \real{0.3784}}
  >{\raggedright\arraybackslash}p{(\linewidth - 2\tabcolsep) * \real{0.6216}}@{}}
\toprule\noalign{}
\begin{minipage}[b]{\linewidth}\raggedright
Archivo
\end{minipage} & \begin{minipage}[b]{\linewidth}\raggedright
Descripción
\end{minipage} \\
\midrule\noalign{}
\endhead
\bottomrule\noalign{}
\endlastfoot
\texttt{001\_initial\_schema.py} & Tablas base: \texttt{organizations}, \texttt{roles}, \texttt{users} \\
\texttt{002\_categories.py} & Tabla \texttt{categories} con jerarquía padre-hijo \\
\texttt{003\_products.py} & Tabla \texttt{products} \\
\texttt{004\_customers.py} & Tabla \texttt{customers} \\
\texttt{005\_inventory\_movements.py} & Tabla \texttt{inventory\_movements} \\
\texttt{006\_sales.py} & Tablas \texttt{sales} y \texttt{sale\_items} \\
\end{longtable}
}

\section{Pruebas}\label{pruebas}

\subsection{Pruebas del backend (pytest)}\label{pruebas-del-backend-pytest}

\begin{Shaded}
\begin{Highlighting}[]
\CommentTok{\# Ejecutar todas las pruebas con cobertura}
\BuiltInTok{cd}\NormalTok{ backend}
\ExtensionTok{uv}\NormalTok{ run bash scripts/test.sh}

\CommentTok{\# Prueba de un archivo específico}
\ExtensionTok{uv}\NormalTok{ run pytest tests/api/routes/test\_sales.py }\AttributeTok{{-}v}

\CommentTok{\# Prueba de una función específica}
\ExtensionTok{uv}\NormalTok{ run pytest tests/api/routes/test\_users.py::test\_get\_users\_superuser\_me }\AttributeTok{{-}v}

\CommentTok{\# Pruebas por patrón}
\ExtensionTok{uv}\NormalTok{ run pytest }\AttributeTok{{-}k} \StringTok{"test\_create"} \AttributeTok{{-}v}

\CommentTok{\# Reporte de cobertura}
\ExtensionTok{uv}\NormalTok{ run coverage run }\AttributeTok{{-}m}\NormalTok{ pytest }\KeywordTok{\&\&} \ExtensionTok{coverage}\NormalTok{ report}
\end{Highlighting}
\end{Shaded}

La cobertura de pruebas abarca los principales módulos del sistema, tanto a nivel de API como de lógica CRUD. Se realizaron pruebas completas de API y CRUD para los módulos de usuarios, categorías, productos, clientes, ventas, movimientos de inventario y dashboard, lo que garantiza la validez funcional y la robustez de estos componentes. En el módulo de autenticación, la cobertura se concentra exclusivamente en pruebas de la API, debido a que la lógica principal reside en los flujos de autenticación y autorización más que en operaciones CRUD tradicionales. Esta cobertura integral asegura que los procesos críticos relacionados con la gestión de datos y la seguridad hayan sido evaluados mediante casos de prueba automatizados.

\subsection{Pruebas del frontend (Playwright E2E)}\label{pruebas-del-frontend-playwright-e2e}

\begin{Shaded}
\begin{Highlighting}[]
\BuiltInTok{cd}\NormalTok{ frontend}

\CommentTok{\# Ejecutar todas las pruebas E2E}
\ExtensionTok{bun}\NormalTok{ run test}

\CommentTok{\# Prueba de un archivo específico}
\ExtensionTok{bunx}\NormalTok{ playwright test tests/login.spec.ts}

\CommentTok{\# Prueba por descripción}
\ExtensionTok{bunx}\NormalTok{ playwright test }\AttributeTok{{-}{-}grep} \StringTok{"crear venta"}

\CommentTok{\# Modo con interfaz visual}
\ExtensionTok{bun}\NormalTok{ run test:ui}
\end{Highlighting}
\end{Shaded}

Las pruebas E2E cubren los flujos: login, registro de organización, registro de usuario, restablecimiento de contraseña, inventario, ventas, clientes, administración de usuarios y configuración.

\section{Generación del Cliente API}\label{generaciuxf3n-del-cliente-api}

Cuando se modifican los endpoints del backend, el cliente TypeScript del frontend debe regenerarse:

\begin{Shaded}
\begin{Highlighting}[]
\CommentTok{\# El backend debe estar corriendo en http://localhost:8000}
\BuiltInTok{cd}\NormalTok{ frontend}
\ExtensionTok{bun}\NormalTok{ run generate{-}client}
\end{Highlighting}
\end{Shaded}

Esto actualiza los archivos en \texttt{src/client/} a partir del esquema OpenAPI del backend. \textbf{No editar manualmente} estos archivos.

\section{Despliegue en Producción}\label{despliegue-en-producciuxf3n}

El despliegue en producción usa el mismo \texttt{compose.yml} sin el \texttt{compose.override.yml}.

\subsection{Prerrequisitos del servidor}\label{prerrequisitos-del-servidor}

\begin{itemize}
\tightlist
\item
  Servidor Linux con Docker Engine y Docker Compose instalados.
\item
  Dominio apuntando a la IP del servidor (registros A/CNAME en el DNS).
\item
  Puertos 80 y 443 abiertos en el firewall.
\end{itemize}

\subsection{Preparar la red de Traefik}\label{preparar-la-red-de-traefik}

Traefik requiere una red Docker externa compartida entre todos los servicios:

\begin{Shaded}
\begin{Highlighting}[]
\ExtensionTok{docker}\NormalTok{ network create traefik{-}public}
\end{Highlighting}
\end{Shaded}

\subsection{Configurar Traefik con TLS}\label{configurar-traefik-con-tls}

Crear el archivo de configuración de Traefik (\texttt{compose.traefik.yml} ya incluido en el repositorio) y levantarlo:

\begin{Shaded}
\begin{Highlighting}[]
\ExtensionTok{docker}\NormalTok{ compose }\AttributeTok{{-}f}\NormalTok{ compose.traefik.yml up }\AttributeTok{{-}d}
\end{Highlighting}
\end{Shaded}

Esto levanta Traefik en modo producción con:

\begin{itemize}
\tightlist
\item
  Redirección automática HTTP → HTTPS.
\item
  Certificados TLS gestionados por Let's Encrypt.
\item
  Panel de Traefik accesible solo internamente.
\end{itemize}

\subsection{Configurar variables de entorno para producción}\label{configurar-variables-de-entorno-para-producciuxf3n}

Editar el \texttt{.env} con los valores de producción:

\begin{Shaded}
\begin{Highlighting}[]
\DataTypeTok{DOMAIN}\OtherTok{=}\StringTok{tu{-}dominio.com}
\DataTypeTok{ENVIRONMENT}\OtherTok{=}\StringTok{production}
\DataTypeTok{SECRET\_KEY}\OtherTok{=}\StringTok{\textless{}clave{-}aleatoria{-}larga\textgreater{}}
\DataTypeTok{POSTGRES\_PASSWORD}\OtherTok{=}\StringTok{\textless{}contraseña{-}segura\textgreater{}}
\DataTypeTok{FIRST\_SUPERUSER\_PASSWORD}\OtherTok{=}\StringTok{\textless{}contraseña{-}segura\textgreater{}}
\DataTypeTok{FRONTEND\_HOST}\OtherTok{=}\StringTok{https://tu{-}dominio.com}
\DataTypeTok{BACKEND\_CORS\_ORIGINS}\OtherTok{=}\StringTok{["https://tu{-}dominio.com"]}
\DataTypeTok{SMTP\_HOST}\OtherTok{=}\StringTok{smtp.tu{-}proveedor.com}
\DataTypeTok{SMTP\_PORT}\OtherTok{=}\DecValTok{587}
\DataTypeTok{SMTP\_TLS}\OtherTok{=}\KeywordTok{true}
\DataTypeTok{EMAILS\_FROM\_EMAIL}\OtherTok{=}\StringTok{noreply@tu{-}dominio.com}
\end{Highlighting}
\end{Shaded}

\subsection{Construir y desplegar}\label{construir-y-desplegar}

\begin{Shaded}
\begin{Highlighting}[]
\CommentTok{\# Construir imágenes con la etiqueta de versión}
\VariableTok{TAG}\OperatorTok{=}\NormalTok{v1.0.0 }\ExtensionTok{docker}\NormalTok{ compose build}

\CommentTok{\# Iniciar todos los servicios}
\VariableTok{TAG}\OperatorTok{=}\NormalTok{v1.0.0 }\ExtensionTok{docker}\NormalTok{ compose up }\AttributeTok{{-}d}

\CommentTok{\# Verificar estado}
\ExtensionTok{docker}\NormalTok{ compose ps}
\ExtensionTok{docker}\NormalTok{ compose logs }\AttributeTok{{-}f}
\end{Highlighting}
\end{Shaded}

Los servicios se levantan en este orden:

\begin{enumerate}
\def\labelenumi{\arabic{enumi}.}
\tightlist
\item
  \texttt{db}: espera estar healthy (hasta 30s).
\item
  \texttt{prestart}: ejecuta migraciones y datos iniciales.
\item
  \texttt{backend} y \texttt{frontend}: arrancan tras \texttt{prestart}.
\item
  \texttt{adminer}: panel de administración de BD.
\end{enumerate}

\subsection{Actualización de la aplicación}\label{actualizaciuxf3n-de-la-aplicaciuxf3n}

\begin{Shaded}
\begin{Highlighting}[]
\CommentTok{\# Obtener la nueva versión}
\FunctionTok{git}\NormalTok{ pull origin main}

\CommentTok{\# Reconstruir imágenes}
\VariableTok{TAG}\OperatorTok{=}\NormalTok{v1.1.0 }\ExtensionTok{docker}\NormalTok{ compose build}

\CommentTok{\# Reiniciar servicios (downtime mínimo: solo el backend se reinicia)}
\VariableTok{TAG}\OperatorTok{=}\NormalTok{v1.1.0 }\ExtensionTok{docker}\NormalTok{ compose up }\AttributeTok{{-}d} \AttributeTok{{-}{-}no{-}deps}\NormalTok{ backend frontend}

\CommentTok{\# Las migraciones se aplican automáticamente al iniciar prestart}
\end{Highlighting}
\end{Shaded}

\section{Monitoreo y Mantenimiento}\label{monitoreo-y-mantenimiento}

\subsection{Estado de los servicios}\label{estado-de-los-servicios}

\begin{Shaded}
\begin{Highlighting}[]
\ExtensionTok{docker}\NormalTok{ compose ps                    }\CommentTok{\# Estado de todos los contenedores}
\ExtensionTok{docker}\NormalTok{ stats                         }\CommentTok{\# Uso de CPU/RAM en tiempo real}
\ExtensionTok{docker}\NormalTok{ compose exec db psql }\AttributeTok{{-}U}\NormalTok{ postgres }\AttributeTok{{-}c} \StringTok{"\textbackslash{}l"}   \CommentTok{\# Listar BDs en PostgreSQL}
\end{Highlighting}
\end{Shaded}

\subsection{Backup de la base de datos}\label{backup-de-la-base-de-datos}

\begin{Shaded}
\begin{Highlighting}[]
\CommentTok{\# Crear un backup}
\ExtensionTok{docker}\NormalTok{ compose exec db pg\_dump }\AttributeTok{{-}U}\NormalTok{ postgres app }\OperatorTok{\textgreater{}}\NormalTok{ backup\_}\VariableTok{$(}\FunctionTok{date}\NormalTok{ +\%Y\%m\%d}\VariableTok{)}\NormalTok{.sql}

\CommentTok{\# Restaurar un backup}
\ExtensionTok{docker}\NormalTok{ compose exec }\AttributeTok{{-}T}\NormalTok{ db psql }\AttributeTok{{-}U}\NormalTok{ postgres app }\OperatorTok{\textless{}}\NormalTok{ backup\_20260401.sql}
\end{Highlighting}
\end{Shaded}

\subsection{Limpiar recursos no usados}\label{limpiar-recursos-no-usados}

\begin{Shaded}
\begin{Highlighting}[]
\ExtensionTok{docker}\NormalTok{ system prune }\AttributeTok{{-}f}              \CommentTok{\# Elimina imágenes, contenedores y redes sin usar}
\ExtensionTok{docker}\NormalTok{ volume prune }\AttributeTok{{-}f}              \CommentTok{\# Elimina volúmenes sin usar (¡precaución en producción!)}
\end{Highlighting}
\end{Shaded}

\section{Solución de Problemas}\label{soluciuxf3n-de-problemas}

{\def\LTcaptype{none} % do not increment counter
\begin{longtable}[]{@{}
  >{\raggedright\arraybackslash}p{(\linewidth - 4\tabcolsep) * \real{0.2044}}
  >{\raggedright\arraybackslash}p{(\linewidth - 4\tabcolsep) * \real{0.1989}}
  >{\raggedright\arraybackslash}p{(\linewidth - 4\tabcolsep) * \real{0.5967}}@{}}
\toprule\noalign{}
\begin{minipage}[b]{\linewidth}\raggedright
Problema
\end{minipage} & \begin{minipage}[b]{\linewidth}\raggedright
Causa probable
\end{minipage} & \begin{minipage}[b]{\linewidth}\raggedright
Solución
\end{minipage} \\
\midrule\noalign{}
\endhead
\bottomrule\noalign{}
\endlastfoot
El backend no inicia & BD no está disponible & Verificar que \texttt{db} esté \texttt{healthy} con \texttt{docker\ compose\ ps}. Revisar \texttt{POSTGRES\_*} en \texttt{.env}. \\
Error de migraciones al iniciar & Migración incompatible & Ejecutar \texttt{docker\ compose\ exec\ backend\ alembic\ history} y \texttt{alembic\ current}. Corregir el script de migración. \\
\texttt{502\ Bad\ Gateway} en el frontend & Backend no responde & Verificar logs del backend: \texttt{docker\ compose\ logs\ backend}. Comprobar el healthcheck. \\
Emails no se envían (producción) & SMTP mal configurado & Verificar \texttt{SMTP\_HOST}, \texttt{SMTP\_PORT} y credenciales. Probar con \texttt{telnet\ \$SMTP\_HOST\ \$SMTP\_PORT}. \\
\texttt{Secret\ key\ too\ short} al iniciar & \texttt{SECRET\_KEY} es demasiado corta & Generar una clave con \texttt{openssl\ rand\ -hex\ 32} y actualizar \texttt{.env}. \\
Traefik no renueva el certificado TLS & Puerto 80 bloqueado o DNS incorrecto & Verificar que el dominio resuelva a la IP correcta y el puerto 80 esté abierto. \\
Error \texttt{UNIQUE\ constraint} al insertar & Dato duplicado en BD & Verificar que el SKU, slug u otro campo único no esté repetido. \\
\end{longtable}
}

\emph{Documento generado como parte del proyecto de grado. Universidad Sergio Arboleda, Semillero de Software como Innovación, Abril 2026.}

    \chapter{Anexo C. Documentación Técnica}\label{anexo-c.-documentaciuxf3n-tuxe9cnica}

\textbf{OrbitEngine}: Plataforma SaaS para la Gestión de Procesos Internos en Pymes\\
Versión: 1.0 \textbar{} Abril 2026\\
Universidad Sergio Arboleda, Semillero de Software como Innovación

\begin{center}\rule{0.5\linewidth}{0.5pt}\end{center}

\section{Stack Tecnológico}\label{stack-tecnoluxf3gico}

\subsection{Backend}\label{backend}

{\def\LTcaptype{none} % do not increment counter
\begin{longtable}[]{@{}
  >{\raggedright\arraybackslash}p{(\linewidth - 6\tabcolsep) * \real{0.2283}}
  >{\raggedright\arraybackslash}p{(\linewidth - 6\tabcolsep) * \real{0.1848}}
  >{\raggedright\arraybackslash}p{(\linewidth - 6\tabcolsep) * \real{0.0761}}
  >{\raggedright\arraybackslash}p{(\linewidth - 6\tabcolsep) * \real{0.5109}}@{}}
\toprule\noalign{}
\begin{minipage}[b]{\linewidth}\raggedright
Componente
\end{minipage} & \begin{minipage}[b]{\linewidth}\raggedright
Tecnología
\end{minipage} & \begin{minipage}[b]{\linewidth}\raggedright
Versión
\end{minipage} & \begin{minipage}[b]{\linewidth}\raggedright
Propósito
\end{minipage} \\
\midrule\noalign{}
\endhead
\bottomrule\noalign{}
\endlastfoot
Lenguaje & Python & 3.10+ & Lenguaje principal del backend \\
Framework web & FastAPI & 0.115+ & API REST asíncrona con tipado estricto \\
ORM + Schemas & SQLModel & 0.0.21+ & Modelos compartidos entre SQLAlchemy y Pydantic \\
Motor de BD & SQLAlchemy & 2.x & Capa de abstracción de base de datos \\
Driver PostgreSQL & Psycopg3 & 3.x & Conector PostgreSQL para Python \\
Migraciones & Alembic & 1.x & Gestión de migraciones de esquema \\
Autenticación & PyJWT & 2.x & Generación y verificación de tokens JWT \\
Hash de contraseñas & Pwdlib & 0.2+ & Hashing seguro con bcrypt \\
Validación & Pydantic & 2.x & Schemas de entrada/salida de datos \\
Configuración & Pydantic-Settings & 2.x & Variables de entorno tipadas \\
Linter/Formatter & Ruff & 0.6+ & Análisis estático y formato de código \\
Verificación de tipos & Mypy & 1.x & Chequeo de tipos estático en modo estricto \\
Testing & Pytest + Coverage & 8.x & Pruebas unitarias e integración \\
Monitoring & Sentry (opcional) & N/A & Rastreo de errores en producción \\
\end{longtable}
}

\subsection{Frontend}\label{frontend}

{\def\LTcaptype{none} % do not increment counter
\begin{longtable}[]{@{}
  >{\raggedright\arraybackslash}p{(\linewidth - 6\tabcolsep) * \real{0.1839}}
  >{\raggedright\arraybackslash}p{(\linewidth - 6\tabcolsep) * \real{0.1839}}
  >{\raggedright\arraybackslash}p{(\linewidth - 6\tabcolsep) * \real{0.0805}}
  >{\raggedright\arraybackslash}p{(\linewidth - 6\tabcolsep) * \real{0.5517}}@{}}
\toprule\noalign{}
\begin{minipage}[b]{\linewidth}\raggedright
Componente
\end{minipage} & \begin{minipage}[b]{\linewidth}\raggedright
Tecnología
\end{minipage} & \begin{minipage}[b]{\linewidth}\raggedright
Versión
\end{minipage} & \begin{minipage}[b]{\linewidth}\raggedright
Propósito
\end{minipage} \\
\midrule\noalign{}
\endhead
\bottomrule\noalign{}
\endlastfoot
Lenguaje & TypeScript & 5.x & Tipado estático sobre JavaScript \\
Framework UI & React & 19.x & Librería de componentes reactivos \\
Build tool & Vite & 6.x & Empaquetado y servidor de desarrollo \\
Router & TanStack Router & 1.x & Enrutamiento con tipado por archivo \\
Server state & TanStack Query & 5.x & Cache y sincronización de datos del servidor \\
Tablas & TanStack Table & 8.x & Tablas con sorting/filtering/paginación \\
Formularios & React Hook Form & 7.x & Manejo de formularios con rendimiento optimizado \\
Validación & Zod & 3.x & Schemas de validación con inferencia de tipos \\
UI Components & Shadcn/ui & N/A & Componentes base sobre Radix UI \\
Estilos & Tailwind CSS & 4.x & Utilidades CSS de bajo nivel \\
Gráficos & Recharts & 2.x & Gráficas para el dashboard \\
HTTP Client & Axios (generado) & N/A & Cliente OpenAPI auto-generado \\
Notificaciones & Sonner & 1.x & Toasts y notificaciones \\
Íconos & Lucide React & 0.4+ & Librería de íconos SVG \\
Linter/Formatter & Biome & 1.x & Análisis estático y formato \\
E2E Testing & Playwright & 1.x & Pruebas end-to-end en navegador \\
\end{longtable}
}

\subsection{Infraestructura}\label{infraestructura}

{\def\LTcaptype{none} % do not increment counter
\begin{longtable}[]{@{}
  >{\raggedright\arraybackslash}p{(\linewidth - 4\tabcolsep) * \real{0.1806}}
  >{\raggedright\arraybackslash}p{(\linewidth - 4\tabcolsep) * \real{0.2361}}
  >{\raggedright\arraybackslash}p{(\linewidth - 4\tabcolsep) * \real{0.5833}}@{}}
\toprule\noalign{}
\begin{minipage}[b]{\linewidth}\raggedright
Componente
\end{minipage} & \begin{minipage}[b]{\linewidth}\raggedright
Tecnología
\end{minipage} & \begin{minipage}[b]{\linewidth}\raggedright
Propósito
\end{minipage} \\
\midrule\noalign{}
\endhead
\bottomrule\noalign{}
\endlastfoot
Base de datos & PostgreSQL 18 & Almacenamiento relacional principal \\
Contenedores & Docker Compose v2 & Orquestación de servicios \\
Reverse proxy & Traefik 3.6 & Enrutamiento HTTP/HTTPS + TLS automático \\
CI/CD & GitHub Actions & Pipeline de integración y entrega continua \\
Panel de BD & Adminer & Administración web de PostgreSQL \\
Email (dev) & Mailcatcher & Interceptor SMTP para desarrollo \\
\end{longtable}
}

\section{Estructura del Proyecto}\label{estructura-del-proyecto}

\begin{verbatim}
orbit-engine/
├── .env                        # Variables de entorno (no incluir en git)
├── compose.yml                 # Configuración Docker Compose base
├── compose.override.yml        # Sobreescrituras para desarrollo local
├── compose.traefik.yml         # Traefik para producción
├── pyproject.toml              # Configuración Ruff y Mypy (nivel raíz)
│
├── backend/
│   ├── Dockerfile
│   ├── pyproject.toml          # Dependencias Python (uv)
│   ├── scripts/
│   │   ├── prestart.sh         # Ejecutado por el servicio prestart
│   │   ├── test.sh             # Script de pruebas con cobertura
│   │   ├── lint.sh             # Mypy + Ruff check
│   │   └── format.sh           # Ruff format
│   └── app/
│       ├── main.py             # Punto de entrada FastAPI + registro de routers
│       ├── models.py           # Todos los modelos SQLModel + schemas Pydantic
│       ├── crud.py             # Operaciones de base de datos
│       ├── utils.py            # Helpers (email, generación de tokens)
│       ├── backend_pre_start.py # Script de inicialización
│       ├── initial_data.py     # Seed de datos iniciales (roles, superusuario)
│       ├── api/
│       │   ├── deps.py         # Dependencias FastAPI (auth, sesión, roles)
│       │   └── routes/
│       │       ├── login.py            # Autenticación y recuperación de contraseña
│       │       ├── users.py            # Gestión de usuarios
│       │       ├── organizations.py    # Registro y configuración de organización
│       │       ├── roles.py            # Listado de roles
│       │       ├── categories.py       # CRUD de categorías
│       │       ├── products.py         # CRUD de productos + stock
│       │       ├── customers.py        # CRUD de clientes
│       │       ├── sales.py            # Registro y gestión de ventas
│       │       ├── inventory_movements.py  # Historial de movimientos
│       │       ├── dashboard.py        # Estadísticas del dashboard
│       │       ├── utils.py            # Health check
│       │       └── private.py          # Endpoints de superusuario
│       ├── core/
│       │   ├── config.py       # Settings con Pydantic-Settings
│       │   ├── db.py           # Engine SQLAlchemy
│       │   └── security.py     # JWT, hashing de contraseñas
│       └── alembic/
│           ├── env.py
│           └── versions/
│               ├── 001_initial_schema.py   # organizations, roles, users
│               ├── 002_categories.py
│               ├── 003_products.py
│               ├── 004_customers.py
│               ├── 005_inventory_movements.py
│               └── 006_sales.py
│
├── frontend/
│   ├── Dockerfile
│   ├── package.json
│   ├── vite.config.ts
│   ├── openapi-ts.config.ts    # Configuración del generador de cliente API
│   └── src/
│       ├── main.tsx            # Punto de entrada React
│       ├── routeTree.gen.ts    # Árbol de rutas auto-generado (no editar)
│       ├── utils.ts            # Helpers compartidos
│       ├── client/             # Cliente API auto-generado (no editar)
│       │   ├── index.ts
│       │   ├── sdk.gen.ts      # Funciones de llamada a la API
│       │   ├── types.gen.ts    # Tipos TypeScript desde OpenAPI
│       │   ├── schemas.gen.ts
│       │   └── core/           # Configuración Axios
│       ├── components/
│       │   ├── Admin/          # Componentes de gestión de usuarios
│       │   ├── Common/         # Componentes reutilizables (DataTable, layouts)
│       │   ├── Customers/      # Componentes del módulo de clientes
│       │   ├── Inventory/      # Componentes del módulo de inventario
│       │   ├── Landing/        # Componentes de la página pública
│       │   ├── Pending/        # Componentes de estados de carga
│       │   ├── Sales/          # Componentes del módulo de ventas
│       │   ├── Sidebar/        # Navegación principal
│       │   ├── UserSettings/   # Configuración de perfil y organización
│       │   ├── theme-provider.tsx
│       │   └── ui/             # Componentes Shadcn/ui (no editar)
│       ├── hooks/
│       │   ├── useAuth.ts
│       │   ├── useCopyToClipboard.ts
│       │   ├── useCustomToast.ts
│       │   ├── useDebounce.ts
│       │   ├── useMobile.ts
│       │   └── useScrollAnimation.ts
│       ├── lib/
│       │   └── utils.ts        # cn() helper para Tailwind
│       └── routes/
│           ├── __root.tsx      # Layout raíz con providers
│           ├── index.tsx       # Landing page (/)
│           ├── login.tsx       # Inicio de sesión (/login)
│           ├── signup.tsx      # Registro de usuario (/signup)
│           ├── signup-org.tsx  # Registro de organización (/signup-org)
│           ├── recover-password.tsx
│           ├── reset-password.tsx
│           └── dashboard/
│               ├── index.tsx   # Dashboard KPIs (/dashboard)
│               ├── inventory.tsx
│               ├── sales.tsx
│               ├── customers.tsx
│               ├── admin.tsx
│               └── settings.tsx
│
└── docs/
    └── informe-final/          # Documentación del proyecto de grado
\end{verbatim}

\section{Autenticación y Autorización}\label{autenticaciuxf3n-y-autorizaciuxf3n}

\subsection{Flujo de autenticación}\label{flujo-de-autenticaciuxf3n}

OrbitEngine usa \textbf{JWT Bearer Tokens} siguiendo el estándar OAuth2 Password Flow:

\begin{enumerate}
\def\labelenumi{\arabic{enumi}.}
\tightlist
\item
  El cliente envía credenciales a \texttt{POST\ /api/v1/login/access-token}.
\item
  El backend valida las credenciales, verifica que la cuenta no esté bloqueada.
\item
  Si son válidas, retorna un \texttt{access\_token} JWT.
\item
  El cliente incluye el token en el header \texttt{Authorization:\ Bearer\ \textless{}token\textgreater{}} en cada request.
\item
  El backend decodifica el token, extrae el \texttt{sub} (UUID del usuario) y carga el usuario de la BD.
\item
  La organización del usuario se extrae de \texttt{user.organization\_id}, no del token directamente.
\end{enumerate}

\textbf{Parámetros del token:}

{\def\LTcaptype{none} % do not increment counter
\begin{longtable}[]{@{}ll@{}}
\toprule\noalign{}
Parámetro & Valor \\
\midrule\noalign{}
\endhead
\bottomrule\noalign{}
\endlastfoot
Algoritmo & HS256 \\
Payload \texttt{sub} & UUID del usuario (\texttt{str}) \\
Expiración & 8 días (11,520 minutos) \\
Clave de firma & \texttt{SECRET\_KEY} del entorno \\
\end{longtable}
}

\subsection{Control de acceso basado en roles (RBAC)}\label{control-de-acceso-basado-en-roles-rbac}

Los roles se almacenan en la tabla \texttt{roles} con permisos como array JSONB. La verificación en los endpoints usa la dependencia \texttt{require\_role()}:

\begin{Shaded}
\begin{Highlighting}[]
\CommentTok{\# Solo administradores pueden crear usuarios}
\AttributeTok{@router.post}\NormalTok{(}\StringTok{"/"}\NormalTok{, dependencies}\OperatorTok{=}\NormalTok{[Depends(require\_role(}\StringTok{"admin"}\NormalTok{))])}
\KeywordTok{def}\NormalTok{ create\_user(...) }\OperatorTok{{-}\textgreater{}}\NormalTok{ Any: ...}

\CommentTok{\# Administradores y vendedores pueden registrar ventas}
\AttributeTok{@router.post}\NormalTok{(}\StringTok{"/"}\NormalTok{, dependencies}\OperatorTok{=}\NormalTok{[Depends(require\_role(}\StringTok{"admin"}\NormalTok{, }\StringTok{"seller"}\NormalTok{))])}
\KeywordTok{def}\NormalTok{ create\_sale(...) }\OperatorTok{{-}\textgreater{}}\NormalTok{ Any: ...}
\end{Highlighting}
\end{Shaded}

\textbf{Dependencias de autenticación disponibles:}

{\def\LTcaptype{none} % do not increment counter
\begin{longtable}[]{@{}
  >{\raggedright\arraybackslash}p{(\linewidth - 2\tabcolsep) * \real{0.2877}}
  >{\raggedright\arraybackslash}p{(\linewidth - 2\tabcolsep) * \real{0.7123}}@{}}
\toprule\noalign{}
\begin{minipage}[b]{\linewidth}\raggedright
Tipo anotado
\end{minipage} & \begin{minipage}[b]{\linewidth}\raggedright
Descripción
\end{minipage} \\
\midrule\noalign{}
\endhead
\bottomrule\noalign{}
\endlastfoot
\texttt{CurrentUser} & Usuario autenticado; error 403 si token inválido \\
\texttt{CurrentAdminUser} & Usuario autenticado con rol \texttt{admin}; error 403 si no \\
\texttt{CurrentOrganization} & UUID de organización extraído del \texttt{CurrentUser} \\
\texttt{SessionDep} & Sesión de base de datos SQLModel \\
\end{longtable}
}

\subsection{Protección contra fuerza bruta}\label{protecciuxf3n-contra-fuerza-bruta}

El modelo \texttt{User} incluye campos \texttt{failed\_login\_attempts} y \texttt{locked\_until}. Tras N intentos fallidos consecutivos, la cuenta se bloquea temporalmente.

\section{API REST. Referencia de Endpoints}\label{api-rest.-referencia-de-endpoints}

Todos los endpoints están prefijados con \texttt{/api/v1}. La documentación interactiva se sirve en \texttt{/docs} (Swagger UI) y \texttt{/redoc}.

\subsection{\texorpdfstring{Autenticación (\texttt{/login})}{Autenticación (/login)}}\label{autenticaciuxf3n-login}

{\def\LTcaptype{none} % do not increment counter
\begin{longtable}[]{@{}
  >{\raggedright\arraybackslash}p{(\linewidth - 6\tabcolsep) * \real{0.0857}}
  >{\raggedright\arraybackslash}p{(\linewidth - 6\tabcolsep) * \real{0.4000}}
  >{\raggedright\arraybackslash}p{(\linewidth - 6\tabcolsep) * \real{0.0571}}
  >{\raggedright\arraybackslash}p{(\linewidth - 6\tabcolsep) * \real{0.4571}}@{}}
\toprule\noalign{}
\begin{minipage}[b]{\linewidth}\raggedright
Método
\end{minipage} & \begin{minipage}[b]{\linewidth}\raggedright
Ruta
\end{minipage} & \begin{minipage}[b]{\linewidth}\raggedright
Auth
\end{minipage} & \begin{minipage}[b]{\linewidth}\raggedright
Descripción
\end{minipage} \\
\midrule\noalign{}
\endhead
\bottomrule\noalign{}
\endlastfoot
POST & \texttt{/login/access-token} & No & Obtener token JWT (login) \\
POST & \texttt{/login/test-token} & Sí & Verificar validez del token \\
POST & \texttt{/password-recovery/\{email\}} & No & Solicitar enlace de recuperación \\
POST & \texttt{/reset-password} & No & Restablecer contraseña con token \\
\end{longtable}
}

\subsection{\texorpdfstring{Usuarios (\texttt{/users})}{Usuarios (/users)}}\label{usuarios-users}

{\def\LTcaptype{none} % do not increment counter
\begin{longtable}[]{@{}
  >{\raggedright\arraybackslash}p{(\linewidth - 8\tabcolsep) * \real{0.0789}}
  >{\raggedright\arraybackslash}p{(\linewidth - 8\tabcolsep) * \real{0.2632}}
  >{\raggedright\arraybackslash}p{(\linewidth - 8\tabcolsep) * \real{0.0526}}
  >{\raggedright\arraybackslash}p{(\linewidth - 8\tabcolsep) * \real{0.1316}}
  >{\raggedright\arraybackslash}p{(\linewidth - 8\tabcolsep) * \real{0.4737}}@{}}
\toprule\noalign{}
\begin{minipage}[b]{\linewidth}\raggedright
Método
\end{minipage} & \begin{minipage}[b]{\linewidth}\raggedright
Ruta
\end{minipage} & \begin{minipage}[b]{\linewidth}\raggedright
Auth
\end{minipage} & \begin{minipage}[b]{\linewidth}\raggedright
Rol mínimo
\end{minipage} & \begin{minipage}[b]{\linewidth}\raggedright
Descripción
\end{minipage} \\
\midrule\noalign{}
\endhead
\bottomrule\noalign{}
\endlastfoot
GET & \texttt{/users/} & Sí & admin & Listar usuarios de la organización \\
POST & \texttt{/users/} & Sí & admin & Crear usuario en la organización \\
GET & \texttt{/users/me} & Sí & Cualquiera & Obtener perfil propio \\
PATCH & \texttt{/users/me} & Sí & Cualquiera & Actualizar nombre/apellido/teléfono \\
DELETE & \texttt{/users/me} & Sí & Cualquiera & Eliminar cuenta propia (soft delete) \\
PATCH & \texttt{/users/me/password} & Sí & Cualquiera & Cambiar contraseña \\
GET & \texttt{/users/\{user\_id\}} & Sí & admin & Obtener usuario por ID \\
PATCH & \texttt{/users/\{user\_id\}} & Sí & admin & Actualizar usuario \\
DELETE & \texttt{/users/\{user\_id\}} & Sí & admin & Eliminar usuario \\
\end{longtable}
}

\subsection{\texorpdfstring{Organizaciones (\texttt{/organizations})}{Organizaciones (/organizations)}}\label{organizaciones-organizations}

{\def\LTcaptype{none} % do not increment counter
\begin{longtable}[]{@{}
  >{\raggedright\arraybackslash}p{(\linewidth - 8\tabcolsep) * \real{0.0952}}
  >{\raggedright\arraybackslash}p{(\linewidth - 8\tabcolsep) * \real{0.2667}}
  >{\raggedright\arraybackslash}p{(\linewidth - 8\tabcolsep) * \real{0.0571}}
  >{\raggedright\arraybackslash}p{(\linewidth - 8\tabcolsep) * \real{0.1238}}
  >{\raggedright\arraybackslash}p{(\linewidth - 8\tabcolsep) * \real{0.4286}}@{}}
\toprule\noalign{}
\multicolumn{5}{@{}>{\raggedright\arraybackslash}p{(\linewidth - 8\tabcolsep) * \real{0.9714} + 8\tabcolsep}@{}}{%
\begin{minipage}[b]{\linewidth}\raggedright
Método \textbar{} Ruta \textbar{} Auth\textbar{} Rol mínimo \textbar{} Descripción \textbar{}
\end{minipage}} \\
\midrule\noalign{}
\endhead
\bottomrule\noalign{}
\endlastfoot
\multicolumn{3}{@{}>{\raggedright\arraybackslash}p{(\linewidth - 8\tabcolsep) * \real{0.4190} + 4\tabcolsep}}{%
GET \textbar{} /organizations/me \textbar{} Sí \textbar{} Cualquiera} & \multicolumn{2}{>{\raggedright\arraybackslash}p{(\linewidth - 8\tabcolsep) * \real{0.5524} + 2\tabcolsep}@{}}{%
Ver datos de la organización actual \textbar{}} \\
\multicolumn{5}{@{}>{\raggedright\arraybackslash}p{(\linewidth - 8\tabcolsep) * \real{0.9714} + 8\tabcolsep}@{}}{%
PATCH \textbar{} /organizations/me \textbar{} Sí \textbar{} admin \textbar{} Actualizar nombre/slug/descripción \textbar{}} \\
\multicolumn{5}{@{}>{\raggedright\arraybackslash}p{(\linewidth - 8\tabcolsep) * \real{0.9714} + 8\tabcolsep}@{}}{%
POST \textbar{} /organizations/signup \textbar{} No \textbar{} N/A \textbar{} Registro de nueva organización + admin \textbar{}} \\
\end{longtable}
}

\subsection{\texorpdfstring{Roles (\texttt{/roles})}{Roles (/roles)}}\label{roles-roles}

{\def\LTcaptype{none} % do not increment counter
\begin{longtable}[]{@{}
  >{\raggedright\arraybackslash}p{(\linewidth - 6\tabcolsep) * \real{0.0968}}
  >{\raggedright\arraybackslash}p{(\linewidth - 6\tabcolsep) * \real{0.1452}}
  >{\raggedright\arraybackslash}p{(\linewidth - 6\tabcolsep) * \real{0.0645}}
  >{\raggedright\arraybackslash}p{(\linewidth - 6\tabcolsep) * \real{0.6935}}@{}}
\toprule\noalign{}
\begin{minipage}[b]{\linewidth}\raggedright
Método
\end{minipage} & \begin{minipage}[b]{\linewidth}\raggedright
Ruta
\end{minipage} & \begin{minipage}[b]{\linewidth}\raggedright
Auth
\end{minipage} & \begin{minipage}[b]{\linewidth}\raggedright
Descripción
\end{minipage} \\
\midrule\noalign{}
\endhead
\bottomrule\noalign{}
\endlastfoot
GET & \texttt{/roles/} & Sí & Listar roles disponibles en la organización \\
\end{longtable}
}

\subsection{\texorpdfstring{Categorías (\texttt{/categories})}{Categorías (/categories)}}\label{categoruxedas-categories}

{\def\LTcaptype{none} % do not increment counter
\begin{longtable}[]{@{}
  >{\raggedright\arraybackslash}p{(\linewidth - 8\tabcolsep) * \real{0.0690}}
  >{\raggedright\arraybackslash}p{(\linewidth - 8\tabcolsep) * \real{0.2069}}
  >{\raggedright\arraybackslash}p{(\linewidth - 8\tabcolsep) * \real{0.0460}}
  >{\raggedright\arraybackslash}p{(\linewidth - 8\tabcolsep) * \real{0.1149}}
  >{\raggedright\arraybackslash}p{(\linewidth - 8\tabcolsep) * \real{0.5632}}@{}}
\toprule\noalign{}
\begin{minipage}[b]{\linewidth}\raggedright
Método
\end{minipage} & \begin{minipage}[b]{\linewidth}\raggedright
Ruta
\end{minipage} & \begin{minipage}[b]{\linewidth}\raggedright
Auth
\end{minipage} & \begin{minipage}[b]{\linewidth}\raggedright
Rol mínimo
\end{minipage} & \begin{minipage}[b]{\linewidth}\raggedright
Descripción
\end{minipage} \\
\midrule\noalign{}
\endhead
\bottomrule\noalign{}
\endlastfoot
GET & \texttt{/categories/} & Sí & Cualquiera & Listar categorías (soporta \texttt{search}, \texttt{is\_active}) \\
POST & \texttt{/categories/} & Sí & admin & Crear categoría \\
GET & \texttt{/categories/\{id\}} & Sí & Cualquiera & Obtener categoría por ID \\
PATCH & \texttt{/categories/\{id\}} & Sí & admin & Actualizar categoría \\
DELETE & \texttt{/categories/\{id\}} & Sí & admin & Eliminar / desactivar categoría \\
\end{longtable}
}

\subsection{\texorpdfstring{Productos (\texttt{/products})}{Productos (/products)}}\label{productos-products}

{\def\LTcaptype{none} % do not increment counter
\begin{longtable}[]{@{}
  >{\raggedright\arraybackslash}p{(\linewidth - 8\tabcolsep) * \real{0.0714}}
  >{\raggedright\arraybackslash}p{(\linewidth - 8\tabcolsep) * \real{0.2540}}
  >{\raggedright\arraybackslash}p{(\linewidth - 8\tabcolsep) * \real{0.0476}}
  >{\raggedright\arraybackslash}p{(\linewidth - 8\tabcolsep) * \real{0.1032}}
  >{\raggedright\arraybackslash}p{(\linewidth - 8\tabcolsep) * \real{0.5000}}@{}}
\toprule\noalign{}
\begin{minipage}[b]{\linewidth}\raggedright
Método
\end{minipage} & \multicolumn{4}{>{\raggedright\arraybackslash}p{(\linewidth - 8\tabcolsep) * \real{0.9048} + 6\tabcolsep}@{}}{%
\begin{minipage}[b]{\linewidth}\raggedright
Ruta \textbar{} Auth\textbar{} Rol mínimo \textbar{} Descripción \textbar{}
\end{minipage}} \\
\midrule\noalign{}
\endhead
\bottomrule\noalign{}
\endlastfoot
\multicolumn{5}{@{}>{\raggedright\arraybackslash}p{(\linewidth - 8\tabcolsep) * \real{0.9762} + 8\tabcolsep}@{}}{%
GET \textbar{} /products/ \textbar{} Sí \textbar{} Cualquiera \textbar{} Listar productos (search, is\_active, category\_id, sort\_by) \textbar{}} \\
\multicolumn{5}{@{}>{\raggedright\arraybackslash}p{(\linewidth - 8\tabcolsep) * \real{0.9762} + 8\tabcolsep}@{}}{%
POST \textbar{} /products/ \textbar{} Sí \textbar{} admin \textbar{} Crear producto \textbar{}} \\
\multicolumn{2}{@{}>{\raggedright\arraybackslash}p{(\linewidth - 8\tabcolsep) * \real{0.3254} + 2\tabcolsep}}{%
GET \textbar{} /products/\{id\} \textbar{} Sí \textbar{} Cualquiera} & \multicolumn{3}{>{\raggedright\arraybackslash}p{(\linewidth - 8\tabcolsep) * \real{0.6508} + 4\tabcolsep}@{}}{%
Obtener producto por ID \textbar{}} \\
\multicolumn{4}{@{}>{\raggedright\arraybackslash}p{(\linewidth - 8\tabcolsep) * \real{0.4762} + 6\tabcolsep}}{%
PATCH \textbar{} /products/\{id\} \textbar{} Sí \textbar{} admin \textbar{} Actualizar producto} & \\
DELETE & \multicolumn{4}{>{\raggedright\arraybackslash}p{(\linewidth - 8\tabcolsep) * \real{0.9048} + 6\tabcolsep}@{}}{%
/products/\{id\} \textbar{} Sí \textbar{} admin \textbar{} Eliminar / desactivar producto \textbar{}} \\
\multicolumn{5}{@{}>{\raggedright\arraybackslash}p{(\linewidth - 8\tabcolsep) * \real{0.9762} + 8\tabcolsep}@{}}{%
POST \textbar{} /products/\{id\}/stock-adjust \textbar{} Sí \textbar{} admin \textbar{} Ajuste manual de stock \textbar{}} \\
\multicolumn{5}{@{}>{\raggedright\arraybackslash}p{(\linewidth - 8\tabcolsep) * \real{0.9762} + 8\tabcolsep}@{}}{%
GET \textbar{} /products/\{id\}/movements \textbar{} Sí \textbar{} Cualquiera \textbar{} Historial de movimientos del producto \textbar{}} \\
\end{longtable}
}

\subsection{\texorpdfstring{Clientes (\texttt{/customers})}{Clientes (/customers)}}\label{clientes-customers}

{\def\LTcaptype{none} % do not increment counter
\begin{longtable}[]{@{}
  >{\raggedright\arraybackslash}p{(\linewidth - 8\tabcolsep) * \real{0.0732}}
  >{\raggedright\arraybackslash}p{(\linewidth - 8\tabcolsep) * \real{0.2805}}
  >{\raggedright\arraybackslash}p{(\linewidth - 8\tabcolsep) * \real{0.0488}}
  >{\raggedright\arraybackslash}p{(\linewidth - 8\tabcolsep) * \real{0.1220}}
  >{\raggedright\arraybackslash}p{(\linewidth - 8\tabcolsep) * \real{0.4756}}@{}}
\toprule\noalign{}
\begin{minipage}[b]{\linewidth}\raggedright
Método
\end{minipage} & \begin{minipage}[b]{\linewidth}\raggedright
Ruta
\end{minipage} & \begin{minipage}[b]{\linewidth}\raggedright
Auth
\end{minipage} & \begin{minipage}[b]{\linewidth}\raggedright
Rol mínimo
\end{minipage} & \begin{minipage}[b]{\linewidth}\raggedright
Descripción
\end{minipage} \\
\midrule\noalign{}
\endhead
\bottomrule\noalign{}
\endlastfoot
GET & \texttt{/customers/} & Sí & Cualquiera & Listar clientes (\texttt{search}, \texttt{is\_active}) \\
POST & \texttt{/customers/} & Sí & Cualquiera & Registrar cliente \\
GET & \texttt{/customers/\{id\}} & Sí & Cualquiera & Obtener cliente por ID \\
PATCH & \texttt{/customers/\{id\}} & Sí & Cualquiera & Actualizar cliente \\
DELETE & \texttt{/customers/\{id\}} & Sí & admin & Desactivar cliente \\
GET & \texttt{/customers/\{id\}/sales} & Sí & Cualquiera & Historial de compras del cliente \\
\end{longtable}
}

\subsection{\texorpdfstring{Ventas (\texttt{/sales})}{Ventas (/sales)}}\label{ventas-sales}

{\def\LTcaptype{none} % do not increment counter
\begin{longtable}[]{@{}
  >{\raggedright\arraybackslash}p{(\linewidth - 8\tabcolsep) * \real{0.0583}}
  >{\raggedright\arraybackslash}p{(\linewidth - 8\tabcolsep) * \real{0.1942}}
  >{\raggedright\arraybackslash}p{(\linewidth - 8\tabcolsep) * \real{0.0388}}
  >{\raggedright\arraybackslash}p{(\linewidth - 8\tabcolsep) * \real{0.0971}}
  >{\raggedright\arraybackslash}p{(\linewidth - 8\tabcolsep) * \real{0.6117}}@{}}
\toprule\noalign{}
\begin{minipage}[b]{\linewidth}\raggedright
Método
\end{minipage} & \begin{minipage}[b]{\linewidth}\raggedright
Ruta
\end{minipage} & \begin{minipage}[b]{\linewidth}\raggedright
Auth
\end{minipage} & \begin{minipage}[b]{\linewidth}\raggedright
Rol mínimo
\end{minipage} & \begin{minipage}[b]{\linewidth}\raggedright
Descripción
\end{minipage} \\
\midrule\noalign{}
\endhead
\bottomrule\noalign{}
\endlastfoot
GET & \texttt{/sales/} & Sí & Cualquiera & Listar ventas (\texttt{search}, \texttt{status}, \texttt{payment\_method}, \texttt{sort\_by}) \\
POST & \texttt{/sales/} & Sí & Cualquiera & Registrar venta (descuenta stock automáticamente) \\
GET & \texttt{/sales/\{id\}} & Sí & Cualquiera & Obtener venta por ID \\
POST & \texttt{/sales/\{id\}/cancel} & Sí & admin & Cancelar venta (revierte stock) \\
\end{longtable}
}

\subsection{\texorpdfstring{Movimientos de inventario (\texttt{/inventory-movements})}{Movimientos de inventario (/inventory-movements)}}\label{movimientos-de-inventario-inventory-movements}

{\def\LTcaptype{none} % do not increment counter
\begin{longtable}[]{@{}
  >{\raggedright\arraybackslash}p{(\linewidth - 6\tabcolsep) * \real{0.0638}}
  >{\raggedright\arraybackslash}p{(\linewidth - 6\tabcolsep) * \real{0.2447}}
  >{\raggedright\arraybackslash}p{(\linewidth - 6\tabcolsep) * \real{0.0426}}
  >{\raggedright\arraybackslash}p{(\linewidth - 6\tabcolsep) * \real{0.6489}}@{}}
\toprule\noalign{}
\begin{minipage}[b]{\linewidth}\raggedright
Método
\end{minipage} & \begin{minipage}[b]{\linewidth}\raggedright
Ruta
\end{minipage} & \begin{minipage}[b]{\linewidth}\raggedright
Auth
\end{minipage} & \begin{minipage}[b]{\linewidth}\raggedright
Descripción
\end{minipage} \\
\midrule\noalign{}
\endhead
\bottomrule\noalign{}
\endlastfoot
GET & \texttt{/inventory-movements/} & Sí & Listar movimientos (\texttt{product\_id}, \texttt{movement\_type}, \texttt{sort\_by}) \\
\end{longtable}
}

\subsection{\texorpdfstring{Dashboard (\texttt{/dashboard})}{Dashboard (/dashboard)}}\label{dashboard-dashboard}

{\def\LTcaptype{none} % do not increment counter
\begin{longtable}[]{@{}
  >{\raggedright\arraybackslash}p{(\linewidth - 6\tabcolsep) * \real{0.0833}}
  >{\raggedright\arraybackslash}p{(\linewidth - 6\tabcolsep) * \real{0.2500}}
  >{\raggedright\arraybackslash}p{(\linewidth - 6\tabcolsep) * \real{0.0556}}
  >{\raggedright\arraybackslash}p{(\linewidth - 6\tabcolsep) * \real{0.6111}}@{}}
\toprule\noalign{}
\begin{minipage}[b]{\linewidth}\raggedright
Método
\end{minipage} & \begin{minipage}[b]{\linewidth}\raggedright
Ruta
\end{minipage} & \begin{minipage}[b]{\linewidth}\raggedright
Auth
\end{minipage} & \begin{minipage}[b]{\linewidth}\raggedright
Descripción
\end{minipage} \\
\midrule\noalign{}
\endhead
\bottomrule\noalign{}
\endlastfoot
GET & \texttt{/dashboard/stats} & Sí & Estadísticas consolidadas de la organización \\
\end{longtable}
}

\textbf{Respuesta de \texttt{/dashboard/stats}:}

\begin{Shaded}
\begin{Highlighting}[]
\FunctionTok{\{}
  \DataTypeTok{"sales\_today"}\FunctionTok{:} \FunctionTok{\{} \DataTypeTok{"count"}\FunctionTok{:} \DecValTok{5}\FunctionTok{,} \DataTypeTok{"total"}\FunctionTok{:} \StringTok{"250000.00"} \FunctionTok{\},}
  \DataTypeTok{"sales\_month"}\FunctionTok{:} \FunctionTok{\{} \DataTypeTok{"count"}\FunctionTok{:} \DecValTok{42}\FunctionTok{,} \DataTypeTok{"total"}\FunctionTok{:} \StringTok{"1850000.00"} \FunctionTok{\},}
  \DataTypeTok{"low\_stock\_count"}\FunctionTok{:} \DecValTok{3}\FunctionTok{,}
  \DataTypeTok{"average\_ticket"}\FunctionTok{:} \StringTok{"44047.62"}\FunctionTok{,}
  \DataTypeTok{"top\_products"}\FunctionTok{:} \OtherTok{[}
    \FunctionTok{\{}
      \DataTypeTok{"product\_id"}\FunctionTok{:} \StringTok{"..."}\FunctionTok{,}
      \DataTypeTok{"product\_name"}\FunctionTok{:} \StringTok{"Producto A"}\FunctionTok{,}
      \DataTypeTok{"quantity\_sold"}\FunctionTok{:} \DecValTok{15}\FunctionTok{,}
      \DataTypeTok{"revenue"}\FunctionTok{:} \StringTok{"450000.00"}
    \FunctionTok{\}}
  \OtherTok{]}\FunctionTok{,}
  \DataTypeTok{"sales\_by\_day"}\FunctionTok{:} \OtherTok{[}\FunctionTok{\{} \DataTypeTok{"date"}\FunctionTok{:} \StringTok{"2026{-}04{-}01"}\FunctionTok{,} \DataTypeTok{"count"}\FunctionTok{:} \DecValTok{3}\FunctionTok{,} \DataTypeTok{"total"}\FunctionTok{:} \StringTok{"120000.00"} \FunctionTok{\}}\OtherTok{]}
\FunctionTok{\}}
\end{Highlighting}
\end{Shaded}

\subsection{Utilidades}\label{utilidades}

{\def\LTcaptype{none} % do not increment counter
\begin{longtable}[]{@{}
  >{\raggedright\arraybackslash}p{(\linewidth - 6\tabcolsep) * \real{0.0923}}
  >{\raggedright\arraybackslash}p{(\linewidth - 6\tabcolsep) * \real{0.3385}}
  >{\raggedright\arraybackslash}p{(\linewidth - 6\tabcolsep) * \real{0.0615}}
  >{\raggedright\arraybackslash}p{(\linewidth - 6\tabcolsep) * \real{0.5077}}@{}}
\toprule\noalign{}
\begin{minipage}[b]{\linewidth}\raggedright
Método
\end{minipage} & \begin{minipage}[b]{\linewidth}\raggedright
Ruta
\end{minipage} & \begin{minipage}[b]{\linewidth}\raggedright
Auth
\end{minipage} & \begin{minipage}[b]{\linewidth}\raggedright
Descripción
\end{minipage} \\
\midrule\noalign{}
\endhead
\bottomrule\noalign{}
\endlastfoot
GET & \texttt{/utils/health-check/} & No & Health check del servicio backend \\
\end{longtable}
}

\section{Modelos de Datos}\label{modelos-de-datos}

\subsection{Organization}\label{organization}

\begin{verbatim}
organizations
├── id              UUID           PK
├── name            VARCHAR(255)   NOT NULL
├── slug            VARCHAR(100)   UNIQUE, INDEX
├── description     TEXT           NULLABLE
├── logo_url        VARCHAR(500)   NULLABLE
├── is_active       BOOLEAN        DEFAULT true
├── created_at      TIMESTAMPTZ    NOT NULL
└── updated_at      TIMESTAMPTZ    NOT NULL

Índices: idx_organizations_slug (unique), idx_organizations_is_active
\end{verbatim}

\textbf{Validación del slug:} expresión regular \texttt{\^{}{[}a-z0-9-{]}\{3,50\}\$} (solo minúsculas, números y guiones).

\subsection{Role}\label{role}

\begin{verbatim}
roles
├── id              INTEGER        PK (autoincrement)
├── name            VARCHAR(50)    UNIQUE, INDEX
├── description     TEXT           NULLABLE
├── permissions     JSONB          DEFAULT []
└── created_at      TIMESTAMPTZ    NOT NULL

Índices: idx_roles_name (unique)
\end{verbatim}

\subsection{User}\label{user}

\begin{verbatim}
users
├── id                     UUID           PK
├── organization_id        UUID           FK → organizations.id, INDEX
├── role_id                INTEGER        FK → roles.id, INDEX
├── email                  VARCHAR(255)   INDEX
├── first_name             VARCHAR(100)   NOT NULL
├── last_name              VARCHAR(100)   NOT NULL
├── phone                  VARCHAR(50)    NULLABLE
├── avatar_url             VARCHAR(500)   NULLABLE
├── hashed_password        VARCHAR(255)   NOT NULL
├── is_active              BOOLEAN        DEFAULT true
├── is_verified            BOOLEAN        DEFAULT false
├── last_login_at          TIMESTAMPTZ    NULLABLE
├── failed_login_attempts  INTEGER        DEFAULT 0
├── locked_until           TIMESTAMPTZ    NULLABLE
├── created_at             TIMESTAMPTZ    NOT NULL
├── updated_at             TIMESTAMPTZ    NOT NULL
└── deleted_at             TIMESTAMPTZ    NULLABLE (soft delete)

Restricciones: UNIQUE(organization_id, email)
Índices: idx_users_organization_id, idx_users_role_id, idx_users_is_active
\end{verbatim}

\subsection{Category}\label{category}

\begin{verbatim}
categories
├── id              UUID           PK
├── organization_id UUID           FK → organizations.id, INDEX
├── parent_id       UUID           FK → categories.id, NULLABLE (jerarquía)
├── name            VARCHAR(255)   NOT NULL
├── description     TEXT           NULLABLE
├── is_active       BOOLEAN        DEFAULT true
├── created_at      TIMESTAMPTZ    NOT NULL
├── updated_at      TIMESTAMPTZ    NOT NULL
└── deleted_at      TIMESTAMPTZ    NULLABLE (soft delete)

Restricciones: UNIQUE(organization_id, name, parent_id)
Índices: idx_categories_organization_id, idx_categories_parent_id, idx_categories_is_active
\end{verbatim}

\subsection{Product}\label{product}

\begin{verbatim}
products
├── id              UUID               PK
├── organization_id UUID               FK → organizations.id, INDEX
├── category_id     UUID               FK → categories.id, NULLABLE
├── name            VARCHAR(255)       NOT NULL
├── sku             VARCHAR(100)       NOT NULL
├── description     TEXT               NULLABLE
├── image_url       VARCHAR(500)       NULLABLE
├── cost_price      NUMERIC(12,2)      DEFAULT 0
├── sale_price      NUMERIC(12,2)      DEFAULT 0
├── stock_quantity  INTEGER            DEFAULT 0
├── stock_min       INTEGER            DEFAULT 0
├── stock_max       INTEGER            NULLABLE
├── unit            VARCHAR(50)        DEFAULT 'unit'
├── barcode         VARCHAR(100)       NULLABLE
├── is_active       BOOLEAN            DEFAULT true
├── created_at      TIMESTAMPTZ        NOT NULL
├── updated_at      TIMESTAMPTZ        NOT NULL
└── deleted_at      TIMESTAMPTZ        NULLABLE (soft delete)

Restricciones: UNIQUE(organization_id, sku)
Índices: idx_products_organization_id, idx_products_category_id, idx_products_is_active, idx_products_sku
\end{verbatim}

\textbf{Stock bajo:} un producto está en alerta cuando \texttt{stock\_quantity\ ≤\ stock\_min}.

\subsection{Customer}\label{customer}

\begin{verbatim}
customers
├── id                UUID           PK
├── organization_id   UUID           FK → organizations.id, INDEX
├── document_type     VARCHAR(50)    NOT NULL (CC, NIT, Pasaporte, Otro)
├── document_number   VARCHAR(50)    NOT NULL
├── first_name        VARCHAR(100)   NOT NULL
├── last_name         VARCHAR(100)   NOT NULL
├── email             VARCHAR(255)   NULLABLE, INDEX
├── phone             VARCHAR(50)    NULLABLE, INDEX
├── address           TEXT           NULLABLE
├── city              VARCHAR(100)   NULLABLE
├── country           VARCHAR(100)   NULLABLE
├── notes             TEXT           NULLABLE
├── is_active         BOOLEAN        DEFAULT true
├── total_purchases   NUMERIC(12,2)  DEFAULT 0 (actualizado automáticamente)
├── purchases_count   INTEGER        DEFAULT 0 (actualizado automáticamente)
├── last_purchase_at  TIMESTAMPTZ    NULLABLE (actualizado automáticamente)
├── created_at        TIMESTAMPTZ    NOT NULL
├── updated_at        TIMESTAMPTZ    NOT NULL
└── deleted_at        TIMESTAMPTZ    NULLABLE (soft delete)

Restricciones: UNIQUE(organization_id, document_number)
Índices: idx_customers_organization_id, idx_customers_email, idx_customers_phone, idx_customers_is_active
\end{verbatim}

\subsection{Sale}\label{sale}

\begin{verbatim}
sales
├── id                   UUID           PK
├── organization_id      UUID           FK → organizations.id, INDEX
├── customer_id          UUID           FK → customers.id, NULLABLE
├── user_id              UUID           FK → users.id, INDEX
├── invoice_number       VARCHAR(100)   NOT NULL
├── sale_date            TIMESTAMPTZ    NOT NULL
├── subtotal             NUMERIC(12,2)  DEFAULT 0
├── discount             NUMERIC(12,2)  DEFAULT 0
├── tax                  NUMERIC(12,2)  DEFAULT 0
├── total                NUMERIC(12,2)  DEFAULT 0
├── payment_method       VARCHAR(50)    DEFAULT 'cash' (cash/card/transfer/other)
├── status               VARCHAR(50)    DEFAULT 'completed' (completed/cancelled/pending)
├── notes                TEXT           NULLABLE
├── cancelled_at         TIMESTAMPTZ    NULLABLE
├── cancelled_by         UUID           FK → users.id, NULLABLE
├── cancellation_reason  TEXT           NULLABLE
├── created_at           TIMESTAMPTZ    NOT NULL
└── updated_at           TIMESTAMPTZ    NOT NULL

Restricciones: UNIQUE(organization_id, invoice_number)
Índices: idx_sales_organization_id, idx_sales_customer_id, idx_sales_user_id,
         idx_sales_status, idx_sales_sale_date, idx_sales_invoice_number
\end{verbatim}

\textbf{Fórmula:} \texttt{total\ =\ subtotal\ -\ discount\ +\ tax}

\subsection{SaleItem}\label{saleitem}

\begin{verbatim}
sale_items
├── id            UUID           PK
├── sale_id       UUID           FK → sales.id, INDEX
├── product_id    UUID           FK → products.id, INDEX
├── product_name  VARCHAR(255)   NOT NULL (snapshot al momento de la venta)
├── product_sku   VARCHAR(100)   NOT NULL (snapshot)
├── quantity      INTEGER        NOT NULL
├── unit_price    NUMERIC(12,2)  NOT NULL
├── subtotal      NUMERIC(12,2)  NOT NULL
└── created_at    TIMESTAMPTZ    NOT NULL

Índices: idx_sale_items_sale_id, idx_sale_items_product_id
\end{verbatim}

\begin{quote}
Los campos \texttt{product\_name} y \texttt{product\_sku} son snapshots del producto en el momento de la venta. Si el producto se edita posteriormente, el historial de ventas conserva los datos originales.
\end{quote}

\subsection{InventoryMovement}\label{inventorymovement}

\begin{verbatim}
inventory_movements
├── id              UUID           PK
├── organization_id UUID           FK → organizations.id, INDEX
├── product_id      UUID           FK → products.id, INDEX
├── user_id         UUID           FK → users.id, INDEX
├── movement_type   VARCHAR(50)    NOT NULL (sale/adjustment/return/purchase)
├── quantity        INTEGER        NOT NULL (positivo = entrada, negativo = salida)
├── previous_stock  INTEGER        NOT NULL
├── new_stock       INTEGER        NOT NULL
├── reference_id    UUID           NULLABLE (FK al objeto de origen)
├── reference_type  VARCHAR(50)    NULLABLE (sale/adjustment/return)
├── reason          TEXT           NULLABLE
└── created_at      TIMESTAMPTZ    NOT NULL

Índices: idx_inventory_movements_organization_id, idx_inventory_movements_product_id,
         idx_inventory_movements_user_id, idx_inventory_movements_movement_type,
         idx_inventory_movements_created_at,
         idx_inventory_movements_reference (reference_type, reference_id)
\end{verbatim}

\section{Variables de Entorno}\label{variables-de-entorno}

La configuración se carga desde el archivo \texttt{.env} usando Pydantic-Settings. Todas las variables se validan al iniciar el servicio.

\subsection{Variables requeridas}\label{variables-requeridas}

{\def\LTcaptype{none} % do not increment counter
\begin{longtable}[]{@{}
  >{\raggedright\arraybackslash}p{(\linewidth - 4\tabcolsep) * \real{0.3662}}
  >{\raggedright\arraybackslash}p{(\linewidth - 4\tabcolsep) * \real{0.1408}}
  >{\raggedright\arraybackslash}p{(\linewidth - 4\tabcolsep) * \real{0.4930}}@{}}
\toprule\noalign{}
\begin{minipage}[b]{\linewidth}\raggedright
Variable
\end{minipage} & \begin{minipage}[b]{\linewidth}\raggedright
Tipo
\end{minipage} & \begin{minipage}[b]{\linewidth}\raggedright
Descripción
\end{minipage} \\
\midrule\noalign{}
\endhead
\bottomrule\noalign{}
\endlastfoot
\texttt{PROJECT\_NAME} & \texttt{str} & Nombre del proyecto \\
\texttt{SECRET\_KEY} & \texttt{str} & Clave para firmar los JWT \\
\texttt{POSTGRES\_SERVER} & \texttt{str} & Host del servidor PostgreSQL \\
\texttt{POSTGRES\_USER} & \texttt{str} & Usuario de PostgreSQL \\
\texttt{POSTGRES\_DB} & \texttt{str} & Nombre de la base de datos \\
\texttt{FIRST\_SUPERUSER} & \texttt{EmailStr} & Email del superusuario inicial \\
\texttt{FIRST\_SUPERUSER\_PASSWORD} & \texttt{str} & Contraseña del superusuario inicial \\
\end{longtable}
}

\subsection{Variables opcionales con valores por defecto}\label{variables-opcionales-con-valores-por-defecto}

{\def\LTcaptype{none} % do not increment counter
\begin{longtable}[]{@{}
  >{\raggedright\arraybackslash}p{(\linewidth - 6\tabcolsep) * \real{0.2353}}
  >{\raggedright\arraybackslash}p{(\linewidth - 6\tabcolsep) * \real{0.1765}}
  >{\raggedright\arraybackslash}p{(\linewidth - 6\tabcolsep) * \real{0.2206}}
  >{\raggedright\arraybackslash}p{(\linewidth - 6\tabcolsep) * \real{0.3529}}@{}}
\toprule\noalign{}
\multicolumn{4}{@{}>{\raggedright\arraybackslash}p{(\linewidth - 6\tabcolsep) * \real{0.9853} + 6\tabcolsep}@{}}{%
\begin{minipage}[b]{\linewidth}\raggedright
Variable \textbar{} Tipo \textbar{} Default \textbar{} Descripción \textbar{}
\end{minipage}} \\
\midrule\noalign{}
\endhead
\bottomrule\noalign{}
\endlastfoot
\multicolumn{4}{@{}>{\raggedright\arraybackslash}p{(\linewidth - 6\tabcolsep) * \real{0.9853} + 6\tabcolsep}@{}}{%
POSTGRES\_PORT \textbar{} int \textbar{} 5432 \textbar{} Puerto de PostgreSQL \textbar{}} \\
\multicolumn{4}{@{}>{\raggedright\arraybackslash}p{(\linewidth - 6\tabcolsep) * \real{0.9853} + 6\tabcolsep}@{}}{%
POSTGRES\_PASS \textbar{} str \textbar{} ``\,'' \textbar{} Contraseña de PostgreSQL \textbar{}} \\
\multicolumn{4}{@{}>{\raggedright\arraybackslash}p{(\linewidth - 6\tabcolsep) * \real{0.9853} + 6\tabcolsep}@{}}{%
ACCESS\_TOKEN \textbar{} int \textbar{} 11520 \textbar{} Expiración del JWT (8 días) \textbar{}} \\
\multicolumn{4}{@{}>{\raggedright\arraybackslash}p{(\linewidth - 6\tabcolsep) * \real{0.9853} + 6\tabcolsep}@{}}{%
FRONTEND\_HOST \textbar{} str \textbar{} http://localhost:5173 \textbar{} URL del frontend para CORS \textbar{}} \\
CORS\_ORIGINS \textbar{} list{[}str{]} \textbar{} {[}{]} & \multicolumn{3}{>{\raggedright\arraybackslash}p{(\linewidth - 6\tabcolsep) * \real{0.7500} + 4\tabcolsep}@{}}{%
Orígenes adicionales CORS \textbar{}} \\
\multicolumn{4}{@{}>{\raggedright\arraybackslash}p{(\linewidth - 6\tabcolsep) * \real{0.9853} + 6\tabcolsep}@{}}{%
ENVIRONMENT \textbar{} Literal \textbar{} ``local'' \textbar{} local \textbar{} staging \textbar{} production \textbar{}} \\
\multicolumn{4}{@{}>{\raggedright\arraybackslash}p{(\linewidth - 6\tabcolsep) * \real{0.9853} + 6\tabcolsep}@{}}{%
SMTP\_HOST \textbar{} str \textbar{} None \textbar{} None \textbar{} Host SMTP para envío de emails \textbar{}} \\
\multicolumn{4}{@{}>{\raggedright\arraybackslash}p{(\linewidth - 6\tabcolsep) * \real{0.9853} + 6\tabcolsep}@{}}{%
SMTP\_PORT \textbar{} int \textbar{} 587 \textbar{} Puerto SMTP \textbar{}} \\
\multicolumn{4}{@{}>{\raggedright\arraybackslash}p{(\linewidth - 6\tabcolsep) * \real{0.9853} + 6\tabcolsep}@{}}{%
SMTP\_TLS \textbar{} bool \textbar{} True \textbar{} Usar TLS en SMTP \textbar{}} \\
\multicolumn{4}{@{}>{\raggedright\arraybackslash}p{(\linewidth - 6\tabcolsep) * \real{0.9853} + 6\tabcolsep}@{}}{%
SMTP\_SSL \textbar{} bool \textbar{} False \textbar{} Usar SSL en SMTP \textbar{}} \\
\multicolumn{4}{@{}>{\raggedright\arraybackslash}p{(\linewidth - 6\tabcolsep) * \real{0.9853} + 6\tabcolsep}@{}}{%
SMTP\_USER \textbar{} str \textbar{} None \textbar{} None \textbar{} Usuario SMTP \textbar{}} \\
\multicolumn{4}{@{}>{\raggedright\arraybackslash}p{(\linewidth - 6\tabcolsep) * \real{0.9853} + 6\tabcolsep}@{}}{%
SMTP\_PASSWORD \textbar{} str \textbar{} None \textbar{} None \textbar{} Contraseña SMTP \textbar{}} \\
\multicolumn{4}{@{}>{\raggedright\arraybackslash}p{(\linewidth - 6\tabcolsep) * \real{0.9853} + 6\tabcolsep}@{}}{%
EMAILS\_FROM \textbar{} EmailStr \textbar{} None \textbar{} None \textbar{} Email remitente \textbar{}} \\
\multicolumn{4}{@{}>{\raggedright\arraybackslash}p{(\linewidth - 6\tabcolsep) * \real{0.9853} + 6\tabcolsep}@{}}{%
EMAIL\_RESET\_TOKEN \textbar{} int \textbar{} 48 \textbar{} Expiración del token de reset \textbar{}} \\
\multicolumn{4}{@{}>{\raggedright\arraybackslash}p{(\linewidth - 6\tabcolsep) * \real{0.9853} + 6\tabcolsep}@{}}{%
SENTRY\_DSN \textbar{} HttpUrl \textbar{} None \textbar{} None \textbar{} DSN de Sentry para error tracking \textbar{}} \\
\end{longtable}
}

\subsection{\texorpdfstring{Propiedades calculadas (no en \texttt{.env})}{Propiedades calculadas (no en .env)}}\label{propiedades-calculadas-no-en-.env}

{\def\LTcaptype{none} % do not increment counter
\begin{longtable}[]{@{}
  >{\raggedright\arraybackslash}p{(\linewidth - 2\tabcolsep) * \real{0.2874}}
  >{\raggedright\arraybackslash}p{(\linewidth - 2\tabcolsep) * \real{0.7126}}@{}}
\toprule\noalign{}
\begin{minipage}[b]{\linewidth}\raggedright
Propiedad
\end{minipage} & \begin{minipage}[b]{\linewidth}\raggedright
Descripción
\end{minipage} \\
\midrule\noalign{}
\endhead
\bottomrule\noalign{}
\endlastfoot
\texttt{SQLALCHEMY\_DATABASE\_URI} & URI completa de conexión a PostgreSQL \\
\texttt{all\_cors\_origins} & \texttt{BACKEND\_CORS\_ORIGINS} + \texttt{FRONTEND\_HOST} \\
\texttt{emails\_enabled} & \texttt{True} si \texttt{SMTP\_HOST} y \texttt{EMAILS\_FROM\_EMAIL} están configurados \\
\end{longtable}
}

\section{Convenciones de Código}\label{convenciones-de-cuxf3digo}

\subsection{Backend (Python)}\label{backend-python}

\textbf{Orden de imports:}

\begin{Shaded}
\begin{Highlighting}[]
\CommentTok{\# 1. Biblioteca estándar}
\ImportTok{import}\NormalTok{ uuid}
\ImportTok{from}\NormalTok{ typing }\ImportTok{import}\NormalTok{ Any}
\ImportTok{from}\NormalTok{ decimal }\ImportTok{import}\NormalTok{ Decimal}

\CommentTok{\# 2. Terceros}
\ImportTok{from}\NormalTok{ fastapi }\ImportTok{import}\NormalTok{ APIRouter, HTTPException}
\ImportTok{from}\NormalTok{ sqlmodel }\ImportTok{import}\NormalTok{ SQLModel, Field}

\CommentTok{\# 3. Locales}
\ImportTok{from}\NormalTok{ app.api.deps }\ImportTok{import}\NormalTok{ CurrentUser, SessionDep}
\ImportTok{from}\NormalTok{ app.models }\ImportTok{import}\NormalTok{ ProductCreate, ProductPublic}
\end{Highlighting}
\end{Shaded}

\textbf{Nomenclatura:}

\begin{itemize}
\tightlist
\item
  Funciones y variables: \texttt{snake\_case}
\item
  Clases y modelos: \texttt{PascalCase}
\item
  Modelos de BD: singular (\texttt{Product}, \texttt{Sale})
\item
  Schemas: \texttt{\{Modelo\}Create}, \texttt{\{Modelo\}Update}, \texttt{\{Modelo\}Public}, \texttt{\{Modelo\}sPublic}
\item
  Aliases de tipo: \texttt{PascalCase} con \texttt{Annotated} (\texttt{CurrentUser}, \texttt{SessionDep})
\end{itemize}

\textbf{Anotaciones de tipo:} obligatorias en todas las funciones. Usar \texttt{str\ \textbar{}\ None} (no \texttt{Optional{[}str{]}}).

\textbf{Patrón de rutas:}

\begin{Shaded}
\begin{Highlighting}[]
\AttributeTok{@router.post}\NormalTok{(}\StringTok{"/"}\NormalTok{, response\_model}\OperatorTok{=}\NormalTok{ProductPublic)}
\KeywordTok{def}\NormalTok{ create\_product(}
    \OperatorTok{*}\NormalTok{,}
\NormalTok{    session: SessionDep,}
\NormalTok{    current\_user: CurrentUser,}
\NormalTok{    current\_organization: CurrentOrganization,}
\NormalTok{    product\_in: ProductCreate,}
\NormalTok{) }\OperatorTok{{-}\textgreater{}}\NormalTok{ Any:}
\NormalTok{    product }\OperatorTok{=}\NormalTok{ Product.model\_validate(product\_in, update}\OperatorTok{=}\NormalTok{\{}\StringTok{"organization\_id"}\NormalTok{: current\_organization\})}
\NormalTok{    session.add(product)}
\NormalTok{    session.commit()}
\NormalTok{    session.refresh(product)}
    \ControlFlowTok{return}\NormalTok{ product}
\end{Highlighting}
\end{Shaded}

\textbf{Códigos HTTP usados:}

\begin{itemize}
\tightlist
\item
  \texttt{400}: Solicitud inválida (datos incorrectos)
\item
  \texttt{403}: Sin permisos / credenciales inválidas
\item
  \texttt{404}: Recurso no encontrado
\item
  \texttt{409}: Conflicto (recurso duplicado)
\end{itemize}

\textbf{Multi-tenancy:} \textbf{todas} las consultas de negocio incluyen filtro \texttt{organization\_id\ =\ current\_organization}. El valor se extrae del usuario autenticado, nunca del cuerpo del request.

\subsection{Frontend (TypeScript/React)}\label{frontend-typescriptreact}

\textbf{Orden de imports:}

\begin{Shaded}
\begin{Highlighting}[]
\CommentTok{// 1. Externos}
\ImportTok{import}\NormalTok{ \{ useState \} }\ImportTok{from} \StringTok{"react"}\OperatorTok{;}
\ImportTok{import}\NormalTok{ \{ useMutation}\OperatorTok{,}\NormalTok{ useQueryClient \} }\ImportTok{from} \StringTok{"@tanstack/react{-}query"}\OperatorTok{;}
\ImportTok{import}\NormalTok{ \{ z \} }\ImportTok{from} \StringTok{"zod"}\OperatorTok{;}

\CommentTok{// 2. Internos (alias @/)}
\ImportTok{import}\NormalTok{ \{ ProductsService \} }\ImportTok{from} \StringTok{"@/client"}\OperatorTok{;}
\ImportTok{import}\NormalTok{ \{ DataTable \} }\ImportTok{from} \StringTok{"@/components/Common/DataTable"}\OperatorTok{;}
\ImportTok{import}\NormalTok{ \{ useCustomToast \} }\ImportTok{from} \StringTok{"@/hooks/useCustomToast"}\OperatorTok{;}
\end{Highlighting}
\end{Shaded}

\textbf{Componentes:} solo funcionales, un componente por archivo, nombre de archivo en \texttt{PascalCase}.

\textbf{Patrón de mutación:}

\begin{Shaded}
\begin{Highlighting}[]
\KeywordTok{const}\NormalTok{ mutation }\OperatorTok{=} \FunctionTok{useMutation}\NormalTok{(\{}
\NormalTok{  mutationFn}\OperatorTok{:}\NormalTok{ (data}\OperatorTok{:}\NormalTok{ ProductCreate) }\KeywordTok{=\textgreater{}}
\NormalTok{    ProductsService}\OperatorTok{.}\FunctionTok{createProduct}\NormalTok{(\{ requestBody}\OperatorTok{:}\NormalTok{ data \})}\OperatorTok{,}
\NormalTok{  onSuccess}\OperatorTok{:}\NormalTok{ () }\KeywordTok{=\textgreater{}} \FunctionTok{showSuccessToast}\NormalTok{(}\StringTok{"Producto creado"}\NormalTok{)}\OperatorTok{,}
\NormalTok{  onError}\OperatorTok{:}\NormalTok{ (err) }\KeywordTok{=\textgreater{}} \FunctionTok{handleError}\NormalTok{(err)}\OperatorTok{,}
\NormalTok{  onSettled}\OperatorTok{:}\NormalTok{ () }\KeywordTok{=\textgreater{}}\NormalTok{ queryClient}\OperatorTok{.}\FunctionTok{invalidateQueries}\NormalTok{(\{ queryKey}\OperatorTok{:}\NormalTok{ [}\StringTok{"products"}\NormalTok{] \})}\OperatorTok{,}
\NormalTok{\})}\OperatorTok{;}
\end{Highlighting}
\end{Shaded}

\textbf{Validación de formularios:} React Hook Form + Zod schemas con \texttt{zodResolver}.

\textbf{No editar:}

\begin{itemize}
\tightlist
\item
  \texttt{src/client/**}: auto-generado desde OpenAPI
\item
  \texttt{src/components/ui/**}: componentes Shadcn/ui
\item
  \texttt{src/routeTree.gen.ts}: auto-generado por TanStack Router
\end{itemize}

\section{Pipeline de CI/CD}\label{pipeline-de-cicd}

El repositorio incluye workflows de GitHub Actions:

\subsection{\texorpdfstring{Workflow de pruebas (en cada PR y push a \texttt{main})}{Workflow de pruebas (en cada PR y push a main)}}\label{workflow-de-pruebas-en-cada-pr-y-push-a-main}

\begin{enumerate}
\def\labelenumi{\arabic{enumi}.}
\tightlist
\item
  \textbf{Backend:} Ejecuta \texttt{uv\ run\ bash\ scripts/test.sh}; incluye Mypy, Ruff y Pytest con cobertura.
\item
  \textbf{Frontend:} Ejecuta \texttt{bun\ run\ lint}; Biome linting y verificación de tipos TypeScript.
\end{enumerate}

\subsection{Workflow de despliegue}\label{workflow-de-despliegue}

Al hacer merge a \texttt{main}:

\begin{enumerate}
\def\labelenumi{\arabic{enumi}.}
\tightlist
\item
  Construye las imágenes Docker de backend y frontend.
\item
  Las etiqueta con el hash del commit y con \texttt{latest}.
\item
  Las publica en el registro de imágenes.
\item
  Actualiza los contenedores en el servidor de producción via SSH.
\end{enumerate}

\section{Estructura de Pruebas}\label{estructura-de-pruebas}

\subsection{Backend (pytest)}\label{backend-pytest}

\begin{verbatim}
backend/tests/
├── conftest.py                     # Fixtures globales: db, client, auth headers
├── api/routes/
│   ├── test_login.py               # Autenticación, recuperación de contraseña
│   ├── test_users.py               # CRUD de usuarios, cambio de contraseña
│   ├── test_categories.py          # CRUD de categorías
│   ├── test_products.py            # CRUD de productos, ajuste de stock
│   ├── test_customers.py           # CRUD de clientes, historial de compras
│   ├── test_sales.py               # Registro y cancelación de ventas
│   ├── test_inventory_movements.py # Listado de movimientos
│   ├── test_dashboard.py           # Estadísticas del dashboard
│   └── test_private.py             # Endpoints de superusuario
└── crud/
    ├── test_user.py
    ├── test_category.py
    ├── test_product.py
    ├── test_customer.py
    ├── test_sale.py
    ├── test_inventory_movement.py
    └── test_dashboard.py
\end{verbatim}

\textbf{Fixtures disponibles en conftest.py:}

{\def\LTcaptype{none} % do not increment counter
\begin{longtable}[]{@{}
  >{\raggedright\arraybackslash}p{(\linewidth - 2\tabcolsep) * \real{0.3462}}
  >{\raggedright\arraybackslash}p{(\linewidth - 2\tabcolsep) * \real{0.6538}}@{}}
\toprule\noalign{}
\begin{minipage}[b]{\linewidth}\raggedright
Fixture
\end{minipage} & \begin{minipage}[b]{\linewidth}\raggedright
Descripción
\end{minipage} \\
\midrule\noalign{}
\endhead
\bottomrule\noalign{}
\endlastfoot
\texttt{db} & Sesión de BD para pruebas (con rollback automático) \\
\texttt{client} & \texttt{TestClient} de FastAPI \\
\texttt{superuser\_token\_headers} & Headers con token del superusuario \\
\texttt{normal\_user\_token\_headers} & Headers con token de usuario normal \\
\end{longtable}
}

\textbf{Ejecutar pruebas:}

\begin{Shaded}
\begin{Highlighting}[]
\BuiltInTok{cd}\NormalTok{ backend}
\ExtensionTok{uv}\NormalTok{ run bash scripts/test.sh           }\CommentTok{\# Todos los tests con cobertura}
\ExtensionTok{uv}\NormalTok{ run pytest tests/api/routes/ }\AttributeTok{{-}v}    \CommentTok{\# Solo tests de API}
\ExtensionTok{uv}\NormalTok{ run pytest }\AttributeTok{{-}k} \StringTok{"test\_create"} \AttributeTok{{-}v}     \CommentTok{\# Tests por patrón}
\end{Highlighting}
\end{Shaded}

\subsection{Frontend (Playwright E2E)}\label{frontend-playwright-e2e}

\begin{verbatim}
frontend/tests/
├── auth.setup.ts                   # Configuración de autenticación compartida
├── config.ts                       # Variables de configuración de tests
├── login.spec.ts                   # Flujo de inicio de sesión
├── sign-up.spec.ts                 # Registro de usuario
├── organization-signup.spec.ts     # Registro de organización
├── reset-password.spec.ts          # Recuperación de contraseña
├── inventory.spec.ts               # Módulo de inventario
├── sales.spec.ts                   # Módulo de ventas
├── customers.spec.ts               # Módulo de clientes
├── admin.spec.ts                   # Gestión de usuarios
├── user-settings.spec.ts           # Configuración de cuenta
└── utils/                          # Helpers de pruebas E2E
\end{verbatim}

\section{Generación del Cliente API}\label{generaciuxf3n-del-cliente-api}

El cliente TypeScript del frontend se genera automáticamente desde el esquema OpenAPI del backend:

\begin{Shaded}
\begin{Highlighting}[]
\CommentTok{\# El backend debe estar corriendo en http://localhost:8000}
\BuiltInTok{cd}\NormalTok{ frontend}
\ExtensionTok{bun}\NormalTok{ run generate{-}client}
\end{Highlighting}
\end{Shaded}

La configuración del generador está en \texttt{frontend/openapi-ts.config.ts}. Los archivos generados en \texttt{src/client/} no deben editarse manualmente; cualquier cambio se sobreescribirá en la próxima generación.

\section{Multi-Tenancy}\label{multi-tenancy}

OrbitEngine implementa multi-tenancy mediante \textbf{aislamiento por discriminador en tablas compartidas} (Shared Database, Separate Rows):

\begin{itemize}
\tightlist
\item
  Todas las tablas de negocio (\texttt{categories}, \texttt{products}, \texttt{customers}, \texttt{sales}, \texttt{sale\_items}, \texttt{inventory\_movements}) incluyen la columna \texttt{organization\_id}.
\item
  La \texttt{organization\_id} se indexa en cada tabla para garantizar rendimiento en consultas filtradas.
\item
  Toda consulta de datos de negocio filtra por \texttt{organization\_id\ =\ current\_user.organization\_id}.
\item
  El valor de \texttt{organization\_id} \textbf{nunca} se acepta como parámetro del request; siempre se extrae del token JWT del usuario autenticado.
\item
  Las restricciones de unicidad incluyen \texttt{organization\_id} como parte de la clave compuesta (ej. \texttt{UNIQUE(organization\_id,\ sku)}), permitiendo que el mismo SKU exista en organizaciones distintas.
\end{itemize}

Esta estrategia garantiza aislamiento completo de datos entre tenants sin necesidad de esquemas o bases de datos separadas, facilitando el mantenimiento y la escalabilidad horizontal.

\emph{Documento generado como parte del proyecto de grado. Universidad Sergio Arboleda, Semillero de Software como Innovación, Abril 2026.}


\end{document}
